% !TEX program = pdflatex

%%%%%%%%%%%%%%
% ShimadaDelPezzoZariski.tex
%%%%%%%%%%%%%%

\pdfoutput=1 
\documentclass{math}
\usepackage{hyperref}
%
\hypersetup{
colorlinks=true,
urlcolor=blue,
citecolor=blue}
%\usepackage[all]{xy,xypic}
\usepackage[all]{xy}
\usepackage{amsfonts,amssymb,amsmath,amsgen,amsopn,amsbsy,theorem,graphicx,epsfig}
\usepackage{eufrak,amscd,bezier,latexsym,mathrsfs,enumerate,multirow}
\usepackage[utf8]{inputenc}\usepackage[english]{babel}
\usepackage[dvipsnames]{xcolor}
\usepackage[pagewise]{lineno}
\linenumbers

\yil{}
\vol{}
\fpage{}
\lpage{}
\doi{}

\title{Del Pezzo surfaces of degree one and examples of Zariski multiples }

\author[Ichiro Shimada]{
\textbf{Ichiro SHIMADA$^{1}$\thanks{Correspondence: ichiro-shimada@hiroshima-u.ac.jp}}\\
$^{1}$Department of Mathematics,
Graduate School of Science,
Hiroshima University,\\
1-3-1 Kagamiyama,
Higashi-Hiroshima,
739-8526 JAPAN
\\ ORCID iD: https://orcid.org/0000-0002-9781-6342\\
\\ [1.8em]

\rec{.202}
\acc{.202}
\finv{..202}
}


\amssayisi{2010 {\itshape AMS Mathematics Subject Classification:} 14H50, 14J26
\newline
\vspace{-15mm}
\begin{center}
		 \raisebox{-17ex}[0ex][0ex]{~~ \raisebox{.5ex}[0ex][0ex]{\footnotesize  This work is licensed under a Creative Commons Attribution 4.0 International License.}}
 		    \end{center}}

\usepackage{tikz}
\usepackage{bm}
\usepackage{enumitem}
%\usetikzlibrary{external}
%\tikzexternalize
%%%%%%%%%%%%%%%%%%%%%%%%%%%%%%%%%%%%%%%%%%%%
%
% ShimadaMacros.tex
%
%%%%%%%%%%%%%%%%%%%%%%%%%%%%%%%%%%%%%%%%%%%%

\theoremstyle{plain}
\newtheorem{theorem}{Theorem}[section]
\newtheorem{lemma}[theorem]{Lemma}
\newtheorem{proposition}[theorem]{Proposition}
\newtheorem{corollary}[theorem]{Corollary}

%\theoremstyle{definition}
\newtheorem{definition}[theorem]{Definition}

%\theoremstyle{remark}
\newtheorem{remark}[theorem]{Remark}

\numberwithin{equation}{section}
%\numberwithin{table}{section}
%\numberwithin{figure}{section}
\newcommand{\qed}{\hfill {$\Box$}}


%FONTS

\renewcommand{\AA}{\mathord{\mathbb A}}
\newcommand{\BB}{\mathord{\mathbb B}}
\newcommand{\CC}{\mathord{\mathbb C}}
\newcommand{\DD}{\mathord{\mathbb D}}
\newcommand{\EE}{\mathord{\mathbb E}}
\newcommand{\FF}{\mathord{\mathbb F}}
\newcommand{\GG}{\mathord{\mathbb G}}
\newcommand{\HH}{\mathord{\mathbb H}}
\newcommand{\II}{\mathord{\mathbb I}}
\newcommand{\JJ}{\mathord{\mathbb J}}
\newcommand{\KK}{\mathord{\mathbb K}}
\newcommand{\LL}{\mathord{\mathbb L}}
\newcommand{\MM}{\mathord{\mathbb M}}
\newcommand{\NN}{\mathord{\mathbb N}}
\newcommand{\OO}{\mathord{\mathbb O}}
\newcommand{\PP}{\mathord{\mathbb  P}}
\newcommand{\QQ}{\mathord{\mathbb  Q}}
\newcommand{\RR}{\mathord{\mathbb R}}
\newcommand{\TT}{\mathord{\mathbb T}}
\newcommand{\UU}{\mathord{\mathbb U}}
\newcommand{\VV}{\mathord{\mathbb V}}
\newcommand{\WW}{\mathord{\mathbb W}}
\newcommand{\XX}{\mathord{\mathbb X}}
\newcommand{\YY}{\mathord{\mathbb Y}}
\newcommand{\ZZ}{\mathord{\mathbb Z}}

\newcommand{\AAA}{\mathord{\mathcal A}}
\newcommand{\BBB}{\mathord{\mathcal B}}
\newcommand{\CCC}{\mathord{\mathcal C}}
\newcommand{\DDD}{\mathord{\mathcal D}}
\newcommand{\EEE}{\mathord{\mathcal E}}
\newcommand{\FFF}{\mathord{\mathcal F}}
\newcommand{\GGG}{\mathord{\mathcal G}}
\newcommand{\HHH}{\mathord{\mathcal H}}
\newcommand{\III}{\mathord{\mathcal I}}
\newcommand{\JJJ}{\mathord{\mathcal J}}
\newcommand{\KKK}{\mathord{\mathcal K}}
\newcommand{\LLL}{\mathord{\mathcal L}}
\newcommand{\MMM}{\mathord{\mathcal M}}
\newcommand{\NNN}{\mathord{\mathcal N}}
\newcommand{\OOO}{\mathord{\mathcal O}}
\newcommand{\PPP}{\mathord{\mathcal P}}
\newcommand{\QQQ}{\mathord{\mathcal Q}}
\newcommand{\RRR}{\mathord{\mathcal R}}
\newcommand{\SSS}{\mathord{\mathcal S}}
\newcommand{\TTT}{\mathord{\mathcal T}}
\newcommand{\UUU}{\mathord{\mathcal U}}
\newcommand{\VVV}{\mathord{\mathcal V}}
\newcommand{\WWW}{\mathord{\mathcal W}}
\newcommand{\XXX}{\mathord{\mathcal X}}
\newcommand{\YYY}{\mathord{\mathcal Y}}
\newcommand{\ZZZ}{\mathord{\mathcal Z}}


%ARROWS

\newcommand{\maprightsp}[1]{\; \smash{\mathop{\; \longrightarrow \; }\limits\sp{#1}}\; }
\newcommand{\mapleftsp}[1]{\; \smash{\mathop{\; \longleftarrow \; }\limits\sp{#1}}\; }

\newcommand{\mapdown}{\phantom{\Big\downarrow}\hskip -8pt \downarrow}
\newcommand{\mapdownleft}[1]{\llap{$\vcenter{\hbox{$\scriptstyle#1$\;}}$}\mapdown}

\newcommand{\inj}{\hookrightarrow}
\newcommand{\surj}{\mathbin{\to \hskip -7pt \to}}
\newcommand{\isom}{\xrightarrow{\raise -3pt \hbox{\scriptsize $\sim$} }}

\newcommand{\set}[2]{\{\,{#1}\mid {#2} \,\}}
\newcommand{\bigset}[2]{\left\{\; {#1} \; \left\vert \; {#2} \;  \right.\right \}}

\newcommand{\angs}[1]{\langle {#1}  \rangle}

\newcommand{\wt}{\widetilde}
\newcommand{\ol}{\overline}


\newcommand{\inv}{\sp{-1}}
\newcommand{\dual}{\sp{\vee}}

\newcommand{\sprime}{\sp{\prime}}
\newcommand{\sperp}{\sp{\perp}}

\newcommand{\semidirectproduct}{\rtimes}

\DeclareMathOperator{\Aut}{Aut}
\DeclareMathOperator{\Image}{Image}
\DeclareMathOperator{\Sing}{Sing}
\DeclareMathOperator{\Ker}{Ker}

\newcommand{\PGL}{\mathord{\mathrm{PGL}}}
\newcommand{\OG}{\mathord{\mathrm{O}}}

\newcommand{\id}{\mathord{\mathrm{id}}}

\newcommand{\rmand}{\textrm{and}}
\newcommand{\rmor}{\textrm{or}}

\newcommand{\mystruth}[1]{\phantom{\hbox{\vrule height #1}}}
\newcommand{\mystrutd}[1]{\phantom{\hbox{\vrule depth #1}}}
\newcommand{\mystruthd}[2]{\phantom{\hbox{\vrule  height #1 depth #2}}}

\newcommand{\pione}{\pi_1}

\newcommand{\intf}[1]{\langle #1 \rangle}

%%%%%%%%%%%%%%%%%%%%%%%%%%%%%%%%%%%%%%%%%%%%

\newcommand{\spset}[1]{^{\{#1\}}}

\newcommand{\mtan}[1]{t_{#1}}
\newcommand{\ADE}{\mathrm{ADE}}
\newcommand{\Tac}{\mathrm{Tac\,}}
\newcommand{\Pic}{\mathrm{Pic}}

\newcommand{\mon}{\Phi}
\newcommand{\LLn}[1]{\LLL^{[#1]}}
\newcommand{\bfP}{\mathbf{P}}

\newcommand{\blambda}{\bm{\lambda}}
\newcommand{\bp}{\bm{p}}
\newcommand{\bP}{\bm{P}}
\newcommand{\brho}{\bm{\rho}}

\newcommand{\sfP}{\mathsf{P}}
\newcommand{\sfB}{\mathsf{B}}
\newcommand{\sfQ}{\mathsf{Q}}

\newcommand{\birat}{\dashrightarrow}

\newcommand{\tilQ}{\widetilde{Q}}
\newcommand{\tilB}{\widetilde{B}}

\newcommand{\barW}{\overline{W}}
\newcommand{\barDelta}{\overline{\Delta}}
\newcommand{\barL}{\overline{L}}
\newcommand{\barH}{\overline{H}}

\newcommand{\involB}{i_B}


%%%%%%%%%%%%%%%%%%%%%%%%%%%%%%%%%%%%%%%%%%%%
%
%   END OF ShimadaMacros.tex
%
%%%%%%%%%%%%%%%%%%%%%%%%%%%%%%%%%%%%%%%%%%%%%



\setcounter{page}{1}
\begin{document}

\maketitle


\begin{abstract}
We construct examples of Zariski 
$N$-tuples with large 
$N$ using the monodromy action 
of the Weyl group of type $E_8$ on the set of $240$ lines in a del Pezzo surface of degree one.
%
\keywords{Del Pezzo surface, monodromy, lattice, Weyl group, Zariski multiple}
\end{abstract}


\section{Introduction}
%
We work over the field of complex numbers.
A \emph{plane curve} means a reduced, possibly reducible, projective plane curve.
%
%
\par
%
Zariski~\cite{ZariskiA, ZariskiB} demonstrated that an equisingular family of plane curves may fail to be connected by 
providing an example of a pair of $6$-cuspidal plane sextics
such that  their complements %in the projective plane 
have non-isomorphic fundamental groups.
Artal~Bartolo~\cite{Artal1994}  revisited Zariski's work, 
and defined  
a \emph{Zariski pair} 
as a pair of plane curves
 that have the same combinatorial type of singularities 
 but differ in their embedding topologies. 
 More generally, a collection of 
$N$ plane curves is called a \emph{Zariski  $N$-tuple}
 if any two  in the collection form a Zariski pair.
 \par
Since~\cite{Artal1994}, many Zariski multiples have been constructed 
using a wide variety of methods.
The construction of Zariski multiples %seems to have served as 
has served as 
a good testing ground for various techniques in the study of embedding topology of plane curves.
See, for example,  the survey paper~\cite{TheSurvey}.
\par
In this paper,
we employ the lattice structure of the middle homology group of the double plane 
branching along the plane curve as an invariant to distinguish topological types~\cite{Shimada2010}.
We enumerate the topological types in our Zariski multiples  by a monodromy argument initiated in~\cite{Shimada2003}.
Del Pezzo surfaces of degree~$1$ and the Weyl group of type~$E_8$ play 
an important role in our investigation.
\par
As the main result, we obtain the following:
%
\begin{theorem}\label{thm:main1}
For each  integer $k$ satisfying $1<k<119$,
there exists a Zariski $N(k)$-tuple $\ZZZ_k$
consisting of plane curves of degree $7+2k$,
where 
%
\begin{equation}\label{eq:Nbinom}
N(k)\ge \frac{1}{348364800}\binom{120}{k}.
\end{equation}
%
\end{theorem}
%
See Section~\ref{sec:main} for the precise description 
of the combinatorial type 
of the plane curves in  $\ZZZ_k$, and 
the exact values of $N(k)$.
We have $N(120-k)=N(k)$, and 
the values of $N(k)$ for small $k$ are as follows:
%
%
%gap> Read("Task250610a.txt");
%1 1 
%2 2 
%3 5 
%4 15 
%5 48 
%6 212 
%7 1116 
%8 7388 
%9 56946 
%
%
\begin{equation}\label{eq:Nks}
\begin{array}{c|ccccccccc}
k & 1 & 2 & 3 & 4 & 5 & 6 &7 &8 &9 \\
\hline 
N(k) & 1& 2 & 5 & 15 & 48 & 212 &1116& 7388& 56946\rlap{\;.}
\end{array}
\end{equation}
%
Putting $k=60$,
we obtain an example of a  Zariski $N(60)$-tuple with 
%
\begin{equation*}\label{eq:largeN}
N(60)>2.77\times 10^{26}.
\end{equation*}
%
%
\par
In our previous paper~\cite{Shimada2022},
we constructed  Zariski multiples by
using the monodromy of the Weyl group of type $E_7$
for a family of del Pezzo surfaces of degree $2$.
For example, we obtained a Zariski $105$-tuple 
in which each member is 
a union of a smooth quartic curve and $14$ lines chosen from its $28$ bitangents.
These examples extended the works~\cite{Bannai2016, Bannai2019, Bannai2020}
by Bannai et al.
We also constructed a series of  Zariski $N$-tuple with $N\to \infty$  by 
means of
 $4$-tangent conics of a smooth quartic curve.
\par
It is a natural step to extend this argument to 
the  Weyl group of type $E_8$ and 
del Pezzo surfaces of degree $1$.
A smooth quartic curve  in~\cite{Shimada2022} is now replaced by a $\mtan{3}$-sextic~(Definition~\ref{def:t3sextic}),
and its $28$ bitangents are replaced by $120$ special tangent conics~(Definition~\ref{def:specialtangent}).
The choices of $k$ conics from the $120$ special tangent conics yield
the Zariski $N(k)$-tuple $\ZZZ_k$.
It seems to be an interesting problem 
 to ask what becomes of the $4$-tangent conics of a smooth quartic curve in the context of a $\mtan{3}$-sextic.
 See~Remark~\ref{rem:K3}.
\par
This paper contains two improvements compared with~\cite{Shimada2022}.
One is that we present a  unified method to compute the monodromy group 
for families of del Pezzo surfaces.
In fact, the proof of~\cite[Theorem~3.1]{Shimada2022} on the monodromy of a family of  del Pezzo surfaces of degree  $2$ 
was incomplete, and we give a rectified argument in the present paper.
Another improvement is that 
we use a simpler reasoning to distinguish topological types in our Zariski multiples.
This simplification is based on an observation (Lemma~\ref{lem:AB})  by Artal~Bartolo.
\par
%
In several parts of the proofs, we rely on brute-force computations carried out by a computer. 
For this purpose, we employ \texttt{GAP}\footnote{The GAP~Group.
{GAP} -- {G}roups, {A}lgorithms, and {P}rogramming, Version  4.13.0, 2024.
\url{https://www.gap-system.org}.}. %~\cite{GAP}.
%A detailed computational data is available from ~\cite{WE8compdata}.
Detailed computational data are available from Zenodo\footnote{\url{https://doi.org/10.5281/zenodo.16148577}.},
or the author's webpage\footnote{\url{https://home.hiroshima-u.ac.jp/ichiro-shimada/ComputationData.html}.}.
%
\par
The plan of this paper is as follows.
In Section~\ref{sec:main}, we precisely state our main result in Theorem~\ref{thm:main2}.
We define the combinatorial type %$\sigma_k$ 
of our  Zariski $N(k)$-tuples, 
and give the exact value of $N(k)$.
In Section~\ref{sec:monodromy}, 
we develop a general theory to compute the monodromy group
for families of del Pezzo surfaces. 
In Section~\ref{sec:lines}, 
we study the lines in a del Pezzo surface of degree $1$
more closely.
In Section~\ref{sec:mtan3anddelPezzo}, 
we relate the theory of del Pezzo surfaces to that of singular plane curves,
and deduce some properties of $\mtan{3}$-sextics
from our discussion about  del Pezzo surfaces of degree $1$.
In Section~\ref{sec:embtop}, 
we investigate the embedding topology of $\mtan{3}$-sextics and 
prove Theorem~\ref{thm:main2}.
We conclude this paper by a remark on the work~\cite{Roulleau2022} about $K3$ surfaces obtained as 
double covers of  del Pezzo surfaces of degree $1$.
%
%
% notation
%\subsection*{Notation}
% for TJM
\par
\medskip
{\bf Notation.}
We use the following notation and conventions.
%
%\begin{enumerate}[label={(\arabic*)}]
%\item 
(1) 
For a set $M$ and a non-negative integer $k$,
we denote  by 
\[
M\spset{k}:=\displaystyle{\binom{M}{k}}
\]
 the set of subsets $S\subset M$ of size $|S|=k$.
%
%
%\item 
(2) The orthogonal group $\OG(L)$ of a lattice $L$ acts on $L$ from the right.
%\item 
(3) For a topological space $T$, 
we write $H_i(T)$ and $H^i(T)$ for $H_i(T;\ZZ)$ and $H^i(T;\ZZ)$, respectively.
%
%\end{enumerate}
%
%
%
\section{Main result}\label{sec:main}
%
In Section~\ref{subsec:sigmak},
we define the combinatorial type $\sigma_k$ 
of plane curves in our Zariski $N(k)$-tuple $\ZZZ_k$.
In  Section~\ref{subsec:Nk}, 
we define the number $N(k)$,  and state 
our main result Theorem~\ref{thm:main2}.
%
\subsection{Combinatorial type 
\texorpdfstring{$\sigma_k$}{sigmak}}\label{subsec:sigmak}
%
%%%%%% Chganged on 2026/Apr/01
\begin{definition}\label{def:combtype}
Two plane curves $D$ and $D\sprime$ are said to 
have \emph{the same combinatorial type}
if there exist  tubular neighborhoods $T\subset \PP^2$ of $D$ and 
$T\sprime\subset \PP^2$ of $D\sprime$
such that 
$(\PP^2, T)$ and $(\PP^2, T\sprime)$ are homeomorphic.
\end{definition}
%
See~\cite[Remark 3]{TheSurvey} for another formulation of the notion of combinatorial type.
%See~\cite[Remark 3]{TheSurvey} for the precise definition of the notion of 
%\emph{combinatorial type  of plane curves}.
%%%%%%%
%
\begin{definition}\label{def:mtan3}
A germ  $(C, \mathbf{0})$ of isolated plane curve singularity
is said to be a \emph{$\mtan{m}$-singularity} if 
$C$ consists of $m$ smooth local branches
and each pair of the local  branches has intersection number $2$.
\end{definition}
%
Note that a $\mtan{2}$-singularity is an ordinary tacnode (an $a_3$-singularity).
%
\begin{definition}\label{def:ellA}
Let $C\subset \PP^2$ be a plane curve with a $\mtan{m}$-singularity at $A\in C$.
The common tangent line $\Lambda\subset \PP^2$ to the local branches of $C$ at $A$
is called the \emph{tangent line to $C$ at $A$}.
\end{definition}
%
\begin{definition}\label{def:t3sextic}
A plane curve $C$  of degree $6$ is 
called a  \emph{$\mtan{3}$-sextic}  
if the singular locus of $C$ consists of a single point $A$, 
and  $A\in C$ is a $\mtan{3}$-singular point.
\end{definition}
%
%
Note that a $\mtan{3}$-sextic
 is irreducible.
We fix a point $A\in \PP^2$ and a line $\Lambda\subset \PP^2$
passing through $A$.
%
\begin{definition}
A $\mtan{3}$-sextic 
with the singular point $A$ and the tangent line $\Lambda$ at $A$
is called a $\mtan{3}$-sextic  \emph{in the frame $(A, \Lambda)$}.
\end{definition}
%
Let $C$ be a $\mtan{3}$-sextic in the frame $(A, \Lambda)$.
%
\begin{definition}\label{def:specialtangent}
A smooth conic $\Gamma$ is said to be a \emph{special tangent conic} of $C$
if the following three conditions  hold; 
%
% for TJM
%\begin{itemize}
%\item 
the conic $\Gamma$ passes through  $A$,  %and 
the union $C+\Gamma$ has  $\mtan{4}$-singularity at $A$, and 
 %\item
at every intersection point of $C$ and $\Gamma$ other than $A$,
the intersection multiplicity is even.
 %\end{itemize}
 \end{definition}
 %
 For a special tangent conic $\Gamma$, 
 we put 
 \[
 \Tac(\Gamma):=\Sing(C+\Gamma)\setminus\{A\}.
 \] 
%
\begin{definition}\label{def:generalset}
We say that 
a set $\{\Gamma_1, \dots, \Gamma_{k}\}$  
consisting of $k$  special tangent conics of $C$  
is said to be \emph{``in a general position"} 
if the following four conditions (a)--(d)  hold.
%
% for TJM
%
%\begin{itemize}
%\item 
(a) Any two of $\Gamma_1, \dots, \Gamma_{k}$ have local intersection number $2$ at $A$.
%\item 
(b) Each $ \Tac(\Gamma_i)$ consists of $3$ tacnodes of $C+\Gamma_i$.
%\item 
(c) The sets $\Tac(\Gamma_1), \dots, \Tac(\Gamma_{k})$ 
are disjoint to each other.
% \item
 d) The singular points  of 
 the union $C+\Gamma_1+\dots +\Gamma_{k}$ 
 other than $A$ and the tacnodes in $\Tac(\Gamma_1), \dots, \Tac(\Gamma_{k})$ 
 are ordinary nodes.
 %\end{itemize}
\end{definition}
%
In Section~\ref{subsec:planecurveswithmtan}, we  prove the following:
%
\begin{proposition}\label{prop:t3}
All  $\mtan{3}$-sextics 
 in the frame $(A, \Lambda)$
are parameterized by a Zariski open subset  $\TTT$
of a $15$-dimensional linear subspace of $|\OOO_{\PP^2}(6)|$.
\end{proposition}
%
For a point $t\in \TTT$,
we denote by $C_t$ the corresponding $\mtan{3}$-sextic.
In Section~\ref{subsec:pip}, we  prove the following:
%
\begin{proposition}\label{prop:t3conics}
Let $t$ be a general point of $\TTT$.
Then $C_t$ 
 has exactly $120$  special tangent conics, and they 
are in a general position.
\end{proposition}
%
% for TJM
%For $t\in \TTT$,
We  denote by $G(C_t)$ the set of special tangent conics of $C_t$.
For $s=\{\Gamma_1, \dots, \Gamma_k\}\in G(C_t)\spset{k}$,
we put
\[
D_{t, s}:=C_t+\Lambda+\Gamma_1+\cdots+\Gamma_k.
\]
%
\begin{definition}\label{def:sigmak}
By Proposition~\ref{prop:t3conics},
the combinatorial type
of $D_{t,s}$ does not depend on the choice of 
$t\in \TTT$ and $s \in G(C_t)\spset{k}$,
provided that $t$ is general in $\TTT$.
We denote this combinatorial type 
by $\sigma_k$.
\end{definition}
%
%\begin{remark}
The curve of combinatorial type $\sigma_k$  
is of degree $7+2k$,  and its singularities consist
of one $\mtan{4+k}$-singular point, $3k$ tacnodes, and $k(k-1)$ ordinary nodes.
It should be noted that the information given by 
a combinatorial type 
 includes,  not only the types of singular points, 
but also more detailed data such as which irreducible components 
correspond to which local branch of each singular point.
See~\cite[Remark 3]{TheSurvey}.
%\end{remark}
%
\subsection{Main Theorem}\label{subsec:Nk}
%
Let $\EE_8$ be the root lattice of type $E_8$,
that is, $\EE_8$ is an even unimodular \emph{negative-definite} lattice of rank $8$
generated by vectors of square norm $-2$.
(Note that we adopt the sign convention opposite to the standard one.)
We denote by $\Delta(\EE_8)$ be the set of vectors of square norm $-2$ in $\EE_8$,
which is of size $240$.
%
%
Let $W(\EE_8)$ be the Weyl group of $\EE_8$,
that is, the subgroup of $\OG(\EE_8)$ generated by reflections with respect to vectors
of square norm $-2$. 
In fact, we have an equality $W(\EE_8)=\OG(\EE_8)$.
See~\eqref{eq:rtimes}.
We then  put
\[
\barW:=W(\EE_8)/\{\pm \id\}, 
\qquad 
\barDelta:=\Delta(\EE_8)/\{\pm \id\}.
\]
Then $\barW$ acts on $\barDelta$ and hence on the set 
$\barDelta\spset{k}$ 
of $k$-element subsets of $\barDelta$.
We define 
\[
N(k):=\textrm{the number of $\barW$-orbits in  $\barDelta^{\{k\}}$.}
\]
%
\par
Finally, we define the topological types of plane curves as follows.
%
\begin{definition}\label{def:homeo}
Two plane curves $D$ and $D\sprime$ are said to 
have \emph{the same embedding topology}
if there exists a homeomorphism between 
$(\PP^2, D)$ and $(\PP^2, D\sprime)$.
\end{definition}
%
Then our main result is stated as follows:
%
\begin{theorem}\label{thm:main2}
Let $o$ be a general point of $\TTT$.
We consider the plane curves $D_{o, s}$ of combinatorial type $\sigma_k$,
where $s$ runs through $G(C_o)\spset{k}$.
Then, 
classifying these curves by the embedding topology yields  exactly $N(k)$ classes.
\end{theorem}
%
Since $|\barDelta|=120$ and $|\barW|=348364800$,
we have the inequality~\eqref{eq:Nbinom}.
Therefore Theorem~\ref{thm:main2} implies Theorem~\ref{thm:main1}.
%
%
%%%%%%%%%%%%%%%%%
%
%
\section{Del Pezzo surfaces}\label{sec:monodromy}
%
In this section, we investigate del Pezzo surfaces.
For the classical results about del Pezzo surfaces, 
we refer the reader to~\cite{Demazure1980, Manin1986}. % and~\cite[Chapter III-3]{Kollar1996}.
In Section~\ref{subsec:Pic}, 
we study the Picard lattice of a  del Pezzo surface.
In Section~\ref{subsec:monodromy},
we present a simple method (Corollary~\ref{cor:index})
for studying   the monodromy  
of a family of del Pezzo surfaces.
In Sections~\ref{subsec:cubic}--\ref{subsec:bianticanonical},
we apply this method 
to natural families of del Pezzo surfaces of degree $3$, $2$, and $1$, respectively. 
%
\subsection{Picard lattice}\label{subsec:Pic}
%
Let $d$ be a positive integer $\le 6$. 
We put 
%\[ for TJM
$n:=9-d$.
%\]
Let $X$ be a del Pezzo surface of degree $d$,
that is, a smooth surface 
whose anti-canonical class 
$\alpha_{X}:=[-K_X]$ 
is ample of self-intersection number $d$.
The Picard lattice $\Pic(X)$ of $X$ is of rank $n+1$,
and is canonically isomorphic to $H^2(X)$.  
There exists a birational morphism 
%
%\begin{equation*}\label{eq:beta}
$\beta\colon X\to {\bfP}^{2}$
%\end{equation*}
%
that is a blowing-up at  distinct $n$ points on ${\bfP}^{2}$,
and $\Pic(X)$ has a basis $h, e_1, \dots, e_n$,
where $h$ is the class of the pullback of a line  on ${\bfP}^{2}$, and $e_1, \dots, e_n$
are the classes of exceptional curves.
With respect to this basis, 
the Gram matrix of $\Pic(X)$  is the diagonal matrix
with diagonal entries $1, -1, \dots, -1$, 
and the anti-canonical   class $\alpha_{X}\in \Pic(X)$
is written as 
%
\begin{equation*}\label{eq:aX}
\alpha_{X}=(3, -1, \dots, -1).
\end{equation*}
%
\begin{table}
\[
% Table1.tex
% Made by Read("Task20250409a.txt");
%
\begin{array}{cccrcc}
d & n & \tau_n & |W(R(X))| &  |\Aut(\tau_n)| & |L(X)|  \\ 
\hline
6 & 3 & A_{1}+A_{2} & {2}^{2} \cdot {3} & 2 & 6 \mystruth{14pt} \\ 
5 & 4 & A_{4} & {2}^{3} \cdot {3} \cdot {5} & 2 & 10  \\ 
4 & 5 & D_{5} & {2}^{7} \cdot {3} \cdot {5} & 2 & 16  \\ 
3 & 6 & E_{6} & {2}^{7} \cdot {3}^{4} \cdot {5} & 2 & 27  \\ 
2 & 7 & E_{7} & {2}^{10} \cdot {3}^{4} \cdot {5} \cdot {7} & 1 & 56  \\ 
1 & 8 & E_{8} & {2}^{14} \cdot {3}^{5} \cdot {5}^{2} \cdot {7} & 1 & 240  \\ 
\end{array}
%
%
\]
\caption{Lattice theoretic data}\label{table:Table1}
\end{table}%
%
The orthogonal complement 
\[
R(X):=(\alpha_{X})\sperp
\]
of $\alpha_{X}$ in $\Pic(X)$ 
is a negative-definite root lattice of type $\tau_n$,
where $\tau_n$ is given in Table~\ref{table:Table1}.
Let $W(R(X))$ denote the Weyl group of the lattice $R(X)$.
The order of $W(R(X))$ is obtained  from the $\ADE$-type $\tau_n$.
%See, for example,~\cite[Section~2.11]{Humphreys1990}.
The root lattice $R(X)$
has a basis $r_1, \dots, r_n$ 
consisting of $(-2)$-vectors 
whose dual graph is the ordinary Dynkin diagram of type $\tau_n$.
We have 
%
\begin{equation}\label{eq:rtimes}
\OG(R(X))=W(R(X))\rtimes \Aut(\tau_n),
\end{equation}
%
where $\Aut(\tau_n)$ is  the group of symmetries  of 
the root system  
$\{r_1, \dots, r_n\}$. % in $R(X)$.
We put
\[
\OG(\Pic(X), \alpha_{X}):=\set{g\in \OG(\Pic(X))}{\alpha_{X}^g=\alpha_{X}}.
\]
Then we have a natural  homomorphism
%
\begin{equation}\label{eq:restriction}
\OG(\Pic(X), \alpha_{X}) \to \OG(R(X))
\end{equation}
%
given by the restriction $g\mapsto g|R(X)$.
It is obvious that~\eqref{eq:restriction} is injective.
%
\begin{proposition}\label{prop:imageisW}
The image of the homomorphism~\eqref{eq:restriction} is equal to $W(R(X))$.
\end{proposition}
%
\begin{proof}
We put $A:=\ZZ\alpha_{X}$, $R:=R(X)$, and consider 
their discriminant groups 
%\[
$d_A:=A\dual/A$ and  $d_R:=R\dual/R$,
%\]
where $A\dual$ and $R\dual$ are 
the dual lattices of $A$ and $R$, respectively.
%Both $d_A$ and $d_R$ are isomorphic to $\ZZ/d\ZZ$.
%Note that $A$ and $R$ are primitive in $\Pic(X)$.
The unimodular lattice $\Pic(X)$,
which is a submodule of $A\dual \oplus R\dual$
containing $A\oplus R$,
gives rise to the graph 
\[
\Pic(X)/(A\oplus R)\;\subset\; d_A\times d_R
\] 
of an isomorphism $d_A\cong d_R$.
%See,~for example,~\cite{theNikulin}.
Hence 
an isometry $g\in \OG(R)$ extends to 
an isometry 
of $\Pic(X)$ that acts on $A$ trivially
if and only if 
$g$ acts on $d_R$ trivially.
It is easy to see that a reflection with respect to a $(-2)$-vector acts trivially on the discriminant group.
On the other hand, 
a direct calculation shows that a  non-trivial element of $\Aut(\tau_n)$ (if there exists any) acts
 non-trivially on $d_R$.
 Hence Proposition~\ref{prop:imageisW} follows from~\eqref{eq:rtimes}.
\end{proof}
%
%
A smooth rational curve $l$ on $X$ is called a \emph{line} if $\intf{l, \alpha_{X}}=1$.
Every line has a self-intersection number $-1$, and 
the set of lines in $X$ is identified with 
\[
L(X):=\set{\lambda\in \Pic(X)}{\intf{\lambda, \lambda}=-1,\;\intf{\lambda, \alpha_{X}}=1}.
\]
By abuse of language,
we sometimes call an element of $L(X)$ a line.
We can enumerate all the elements of $L(X)$ explicitly.
Let  $L^{[n]}(X)$ denote the set of all ordered $n$-tuples 
%\[
$\blambda=[\lambda_1, \dots, \lambda_n]$
%\]
of lines such that $\intf{\lambda_i, \lambda_j}=0$ for any $i, j$ with  $i\ne j$.
%Note that we have 
%\[
%\bm{e}:=[e_1, \dots, e_n]\in L^{[n]}(X),
%\]
%where $e_1, \dots, e_n$ are the classes of the exceptional curves of~\eqref{eq:beta}.
By Proposition~\ref{prop:imageisW} and~\cite[Proposition II-4]{Demazure1980},  we have the following:
%
\begin{corollary}\label{cor:freeandtransitive}
The natural action of $\OG(\Pic(X), \alpha_{X})\cong W(R(X))$ 
on $L^{[n]}(X)$ is free and transitive.
\qed
\end{corollary}
%
For $\blambda=[\lambda_1, \dots, \lambda_n]\in L^{[n]}(X)$,
we have a birational morphism
to a projective plane 
that is the contraction of the lines 
$l_1, \dots, l_n$ whose classes are $\lambda_1, \dots, \lambda_n$.
We denote this blowing-down morphism by
%
\begin{equation}\label{eq:betalambda}
\beta_{\blambda}\;\colon\; X \;\to\; \bP(X/\blambda).
\end{equation}
%
%
\subsection{Monodromy}\label{subsec:monodromy}
%
Let $f\colon \XXX \to \UUU$ be a family of del Pezzo surfaces of degree  $d$.
We assume that the parameter space $\UUU$ is smooth and irreducible.
For a point $u$ of $\UUU$, we denote by $X_u$ the fiber $f\inv(u)$.
Instead of $\alpha_{X_u}$, we denote by 
$\alpha_u\in \Pic(X_u)$ 
the anti-canonical class of $X_u$.
We choose a base point $b\in \UUU$.
The local system 
$R^2f_* \ZZ$ over $\UUU$
%\[
%R^2f_* \ZZ\; \to\; \UUU
%\]
is a family of the lattices $\Pic(X_u)\cong H^2(X_u)$,
and it has a section $u\mapsto \alpha_{u}$.
Hence we obtain a monodromy homomorphism 
%
\begin{equation}\label{eq:mon}
\mon\;\colon\; \pione(\UUU, b)\;\to\; \OG(\Pic(X_b), \alpha_b)\cong W(R(X_b)).
\end{equation}
%
We investigate  the surjectivity of the monodromy homomorphism $\mon$. %~\eqref{eq:mon}.
%
\par
%
We consider the \'etale covering
\[
\LLL^{[n]}\to\UUU
\]
whose fiber over $u\in\UUU$ is the set $L^{[n]}(X_u)$ of ordered sets of disjoint $n$ lines in $X_u$.
The associated monodromy action 
of $\pione(\UUU, b)$ on $L^{[n]}(X_b)$ 
is given by $[\gamma]\mapsto(\blambda\mapsto\blambda^{ \mon([\gamma])})$
%\[
%[\gamma]\;\;\mapsto\;\;(\blambda\;\mapsto \; \blambda^{ \mon([\gamma])})
%\]
for $[\gamma]\in \pione(\UUU, b)$,
where $\blambda\mapsto \blambda^{ \mon([\gamma])}$  
denotes the action of the element $ \mon([\gamma])$ of 
$ \OG(\Pic(X_b),  \alpha_b)$ on $L^{[n]}(X_b)$.
The orbits of the monodromy action on $L^{[n]}(X_b)$ 
correspond bijectively to the connected components of $\LLL^{[n]}$.
Hence, by Corollary~\ref{cor:freeandtransitive}, we obtain the following:
%
\begin{corollary}\label{cor:index}
The index of the image %$\Gamma:=\Image \mon$ 
of  the monodromy homomorphism~$\mon$ in $W(R(X_b))$ is 
equal to the number of the connected components of $\LLL^{[n]}$.
\qed
\end{corollary}
%
We fix a projective plane ${\bfP}^{2}$.
For an ordered set  $\bp= [p_1, \dots, p_n]\in ({\bfP}^{2})^n$ of distinct $n$ points of ${\bfP}^{2}$, 
we denote by 
$\beta_{\bp}\colon  Y_{\bp}\to {\bfP}^{2}$
%\[
%\beta_{\bp}\;\;\colon\;\;  Y_{\bp}\;\to\; {\bfP}^{2}
%\]
 the blowing-up at 
the points $p_1, \dots, p_n$.
%
\begin{definition}\label{def:PPPn}
Let $\PPP_n$ be the Zariski open subset of $({\bfP}^{2})^n$
consisting of all ordered sets  $\bp=[p_1, \dots, p_n]$ of distinct $n$ points of ${\bfP}^{2}$
such that 
%
%\begin{enumerate}[label={\rm{(\roman *)}}]
%\item  
(i) no three points in  $\bp$ are on a line,
%\item  
(ii) no six points in  $\bp$ are on a conic, and 
%\item   
(iii) there exists no cubic curve passing through $7$  points in  $\bp$
and having a double point at the $8$th point.
%\end{enumerate}
%
\end{definition}
%
\begin{theorem}[Th\'eor\`eme~II-1 in~\cite{Demazure1980}]
\label{thm:PPPn}
The surface  $Y_{\bp}$ is a del Pezzo surface of degree $d=9-n$ if and only if $\bp\in \PPP_n$.
\qed
\end{theorem}
%
For $\bp= [p_1, \dots, p_n] \in \PPP_n$, we have a distinguished element 
%\[
$\brho_{\bp}=[\rho_1, \dots, \rho_n]$
%\]
of $L^{[n]}(Y_{\bp})$, 
where $\rho_i$ is the class of the exceptional curve over $p_i$.
We consider the incidence  variety
%
\[
\III:=\bigset{(u,\bp, g)}{%
\parbox{8.3cm }{\,$u \in \UUU$, \, $\bp\in \PPP_n$, \,
and $g$ is an isomorphism $X_u\isom Y_{\bp}$
}%
}
\]
%
with the projections
\[
\UUU\;\mapleftsp{\varpi_{1}}\;\III\;\maprightsp{\varpi_2}\; \PPP_n.
\]
A point of $\LLL^{[n]}$ is written as $(u, \blambda)$,
where $u\in \UUU$ and $\blambda\in L^{[n]}(X_u)$.
Then we can lift $\varpi_1\colon \III\to \UUU$  to 
$\varpi_1^{\LLL}\colon \III\to \LLL^{[n]}$
by setting
%\[
$\varpi_1^{\LLL}(u, \bp, g):=(u, \blambda)$,
%\]
where $\blambda$ is the element of 
$L^{[n]}(X_u)$ that is mapped to $\brho_{\bp}\in L^{[n]}(Y_{\bp})$ by 
the bijection $L^{[n]}(X_u) \isom L^{[n]}(Y_{\bp}) $ 
induced by  $g\colon X_u\isom Y_{\bp}$:
\[
\begin{array}{ccccc}
\LLL^{[n]} & \mapleftsp{\varpi_1^{\LLL}} & \III &  \maprightsp{\varpi_2} & \PPP_n\mystruth{15pt} \\
\mapdownleft{} &\rlap{$\swarrow$\scriptsize$\varpi_1$} &&& \\
\,\UUU\rlap{\;.}&&&&
\end{array}
\]
The fiber of $\varpi_1^{\LLL}$ over $(u, \blambda)\in \LLL^{[n]}$ is the variety of all isomorphisms
%\[
%\bP(X_u/\blambda)\;\isom\;{\bfP}^{2},
$\bP(X_u/\blambda)\isom{\bfP}^{2}$,
%\]
where $\bP(X_u/\blambda)$ is defined in~\eqref{eq:betalambda}.
Indeed, if we have $\varpi_1^{\LLL}(u,\bp, g)=(u, \blambda)$, then the isomorphism $g\colon X_u\isom Y_{\bp}$ maps  
the exceptional curves of the blowing-down $\beta_{\blambda}\colon X_u\to  \bP(X_u/\blambda)$ to
the exceptional curves of the blowing-down $\beta_{\bp}\colon Y_{\bp}\to {\bfP}^{2}$, 
and hence 
$g$ induces an isomorphism $\bP(X_u/\blambda)\isom{\bfP}^{2}$.
Conversely, if we are given a point $(u, \blambda)$ of $\LLL^{[n]}$ and an isomorphism $\bar{g}\colon \bP(X_u/\blambda)\isom{\bfP}^{2}$,
then, setting $\bp$ to be the image of the centers of 
$\beta_{\blambda}$ by $\bar{g}$ 
and lifting $\bar{g}$
to the isomorphism $g\colon X_u\isom Y_{\bp}$ between  their  blow-ups, 
we obtain  a point $(u,\bp, g)$  in the fiber of $\varpi_1^{\LLL}$ over $(u, \blambda)$.
\par
The variety of isomorphisms %between the projective planes $\bP(X_u/\blambda)$ and ${\bfP}^{2}$
$\bP(X_u/\blambda)\cong {\bfP}^{2}$
is isomorphic to $\PGL(3, \CC)$.
Hence $\varpi_1^{\LLL}$ gives a bijection from the set of connected components of $\III$
to that of $\LLL^{[n]}$.
We investigate the connected components  of $\III$ by 
the second projection $\varpi_2\colon \III \to \PPP_n$.
%
%
%\subsection{The case $d=3$}
\subsection{Family of cubic surfaces}\label{subsec:cubic}
%
The anti-canonical model of a del Pezzo surface of degree $d=3$
is a smooth cubic surface.
%Conversely, every smooth cubic surface 
%is the anti-canonical model of a del Pezzo surface of degree $3$.
%\par
We fix a projective space $\PP^3$,
and consider the family $\XXX\to \UUU$ of smooth cubic surfaces, where 
$\UUU$ is  the Zariski open subset of 
$|\OOO_{\PP^3} (3)|\cong \PP^{19}$
parameterizing all smooth cubic surfaces.
\par
The following reproduces the result of Harris~\cite{Harris1979}  on the Galois group of $27$ lines in a smooth cubic surface.
%
\begin{proposition}\label{prop:WE6}
For the  family 
 $\XXX\to \UUU$ 
 of smooth cubic surfaces,
 the monodromy homomorphism $\mon$  %in~\eqref{eq:mon}
 is surjective onto the Weyl group $W(R(X_b))$ 
 of type $E_6$.
\end{proposition}
%
\begin{proof}
For any point $\bp\in \PPP_6$,
the fiber of $\varpi_2$ over $\bp$ is 
the variety of all isomorphisms between $\PP^3$
and the projective space 
\[
|\alpha_{\bp}|\dual=\PP^{*}(H^0(Y_{\bp}, \OOO(\alpha_{\bp}))),
\]
where $\alpha_{\bp}$ is the anti-canonical class of $Y_{\bp}$.
This variety is isomorphic to $\PGL(4, \CC)$.
Hence $\LLn{6}$ is connected.
Therefore $\mon$ is surjective by Corollary~\ref{cor:index}.
\end{proof}
%
%%%%%%%%%%%%%%%%%Changed 2026/Apr/01
See~\cite{Shimada2026} for an application of Proposition~\ref{prop:WE6}.
%%%%%%%%%%%%%%%%%
%
%\subsection{The case $d=2$}
\subsection{Family of quartic double planes}\label{subsec:quartic}
%
The anti-canonical model of a del Pezzo surface
$X$ of degree $d=2$ is a double plane
$X\to \PP^2$ branching along a smooth quartic curve.
%Conversely, every double plane 
%branching along a smooth quartic curve
%is the anti-canonical model of a del Pezzo surface of degree $2$.
%%
%\par
%
We fix a projective plane $\PP^2$. 
Let  $\UUU$ be the Zariski open subset of 
$|\OOO_{\PP^2} (4)|\cong \PP^{14}$
that parameterizes all smooth quartic curves in $\PP^2$.
For each $u\in \UUU$,
we denote by  $B_u\subset \PP^2$  the corresponding quartic curve.
We consider the family 
 $\XXX\to \UUU$ 
 of smooth quartic double planes,
that is, 
 $\XXX$ is the double cover of $\PP^2\times \UUU$ with the projection  $ \XXX\to \UUU$ 
 whose fiber  $X_u$ over $u\in \UUU$
 is  the double plane $X_u\to \PP^2$
branching along $B_u$.
%
\par
The following proposition was stated in~\cite{Shimada2022}, but the proof was incomplete.
A weaker result concerning
the Galois group of the $28$ bitangents of a smooth quartic curve had been proved  in  Harris~\cite{Harris1979}.
%
\begin{proposition}\label{prop:WE7}
For the  family 
 $ \XXX\to \UUU$ 
 of smooth quartic double planes,
 the monodromy homomorphism $\mon$ %in~\eqref{eq:mon}
 is surjective onto the Weyl group $W(R(X_b))$ 
 of type~$E_7$.
\end{proposition}
%
%
\begin{proof}
Let $\bp$ be a point of $\PPP_7$,
and let $Y_{\bp} \to {\sfP}_{\bp}^2$ 
be the anti-canonical model of $Y_{\bp}$.
Let $\sfB_{\bp}\subset {\sfP}_{\bp}^2$ be the branch curve of $Y_{\bp} \to {\sfP}_{\bp}^2$.
The fiber of $\varpi_2\colon \III\to \PPP_7$ 
over $\bp$ consists 
of pairs $(u, g)$,
where $u$ is a point of $\UUU$
and $g$ is an isomorphism from $X_u$ to $Y_{\bp}$.
An isomorphism $g$ from $X_u$ to $Y_{\bp}$
%for TJM
induces an isomorphism from $\PP^2$ to ${\sfP}_{\bp}^2$.
%
%\[
%\bar{g}\;\colon \; \PP^2 \isom  {\sfP}_{\bp}^2.
%\]
Conversely,
suppose that an isomorphism $\gamma\colon \PP^2\isom {\sfP}_{\bp}^2$ is given.
Let $u\in \UUU$ be the point such that $B_{u}=\gamma\inv(\sfB_{\bp})$.
Then $\gamma$ admits exactly \emph{two}  liftings
%for TJM
$g_1\colon   X_u \isom Y_{\bp}$
and $g_2\colon  X_u \isom Y_{\bp}$,
%\[
%g_1\;\colon \;  X_u \isom Y_{\bp},
%\quad
%g_2\;\colon \;  X_u \isom Y_{\bp},
%\]
which differ by the deck-transformation of $X_u$ over  $\PP^2$.
Since $\PGL(3, \CC)$ is smooth and irreducible,
we see that the fiber of $\varpi_2\colon \III\to \PPP_7$ 
over $\bp$ has at most two connected components.
Therefore
$\III$  has at most two connected components,
and so does $\LLn{7}$.
Consequently,
the index $[W: \Gamma]$ of the image
%\[
$\Gamma:=\Image(\mon)$
%\]  
of $\mon$ in $W:=W(R(X_b))$  is at most $2$.
\par
We assume 
%
\begin{equation}\label{eq:absurdum}
[W: \Gamma]=2, 
\end{equation}
%
and derive a contradiction.
Note  that the Weyl group $W$ of type $E_7$ contains a simple group $G$  
as the kernel of $\det \colon W \to \{\pm 1\}$.
See~\cite[page 46]{ATLAS1985}.
We show in the next paragraph  that  $\Gamma$  contains an element of determinant $-1$.
By assumption~\eqref{eq:absurdum}, we see that 
$G\cap \Gamma$ is a normal subgroup of $G$ with index $2$,
which contradicts the simplicity of $G$.
\par
Let $\HHH\subset |\OOO_{\PP^2}(4)|\cong \PP^{14}$ 
be the hypersurface that parameterizes singular quartic curves.
Then $\HHH$ is irreducible. 
Let $q$ be a general point of $\HHH$.
Then the corresponding quartic curve $B_q\subset \PP^2$ 
has an  ordinary node as its only singularity, 
and hence the double plane $X_q\to \PP^2$ branching along $B_q$ has 
an ordinary double point as its only singularity.
We choose a sufficiently small closed disc $D\subset |\OOO_{\PP^2}(4)|$ 
intersecting $\HHH$ at $q$ transversely,
and let
$\gamma\colon[0,1]\to \UUU$ be a loop 
that goes from the base point $b$ to a point $q\sprime\in \partial D$ along a path $\tau$,
makes a round trip along $\partial D$ in  positive-direction,
and retraces  the  path $\tau$ backwards to $b$.
The monodromy action on $H^2(X_b)$ by   $[\gamma]\in \pione(\UUU, b)$ is calculated by
the \emph{local Picard--Lefschetz formula}.
(See, for example, \cite{Lamotke1981}.)
We have  a \emph{vanishing cycle} $v\in H^2(X_b)$
corresponding to the ordinary double point of $X_q$,
which satisfies $\intf{ \alpha_b, v} =0$ and $\intf{v, v} =-2$,
and the monodromy  on $H^2(X_b)\cong \Pic(X_b)$ by $[\gamma]$ is the reflection $x\mapsto x+\intf{v, x}v$
with respect to $v$.
Hence we have $\det(\mon([\gamma]))=-1$. 
Note that the identification $\OG(\Pic(X_b), \alpha_b)\cong W(R(X_b))$ 
preserves the determinant.
Therefore $\Gamma$  contains an element of determinant $-1$.
%
%
\end{proof}
%
%
\subsection{Family of bi-anti-canonical models of del Pezzo surfaces of degree
\texorpdfstring{$1$}{one}}\label{subsec:bianticanonical}
%
Let $X$ be a del Pezzo surface of degree $d=1$.
Then 
the complete linear system $|2\alpha_X| $
of bi-anti-canonical divisors of $X$ gives rise to a double covering
\[
X\to Q\subset\PP^3
\]
of a singular quadric surface $Q$ of rank $3$ (a quadric cone)
that branches along $B\cup\{V\}$,
where $B$ is a smooth member of $|\OOO_{Q}(3)|$ and $V\in Q$ is the vertex.
Conversely, 
for a quadric cone $Q$ with the vertex $V\in Q$ and 
a smooth member $B$ of  $|\OOO_{Q}(3)|$,
the double cover $X\to Q$ branching  along $B\cup\{V\}$
is the bi-anti-canonical model of a del Pezzo surface $X$ of degree $1$.
%
\par
%
We fix 
a quadric cone $Q$ with the vertex $V\in Q$.
Let $\UUU$ denote the Zariski open subset of 
$|\OOO_{Q} (3)|\cong \PP^{15}$
that parameterizes all smooth members.
For $u\in \UUU$,
we denote by  $B_u\subset Q$  the corresponding curve.
We consider the family 
 $\XXX\to \UUU$ 
  of del Pezzo surfaces 
of degree $1$
such that
 $\XXX$ is the double cover of $Q\times \UUU$ with  the  projection  $\XXX\to \UUU$ 
 whose fiber  $X_u$ over $u\in \UUU$
 is the double cover $X_u\to Q$
branching along $B_u\cup \{V\}$.
%
\begin{proposition}\label{prop:WE8}
For the  family 
 $ \XXX\to \UUU$ 
 of bi-anti-canonical models of del Pezzo surfaces of degree $1$,
 the monodromy homomorphism $\mon$
 is surjective onto the Weyl group $W(R(X_b))$ 
 of type $E_8$.
\end{proposition}
%
%
\begin{proof}
The first half of the proof is almost the same as that of Proposition~\ref{prop:WE7}.
Let $\bp$ be a point of $\PPP_8$, 
and let $Y_{\bp} \to {\sfQ}_{\bp}$ 
be the bi-anti-canonical model of $Y_{\bp}$.
Let $\sfB_{\bp}\subset {\sfQ}_{\bp}$ be the curve component of the branch locus  of $Y_{\bp} \to {\sfQ}_{\bp}$.
The fiber of $\varpi_2\colon \III\to \PPP_8$ 
over $\bp$ consists 
of pairs $(u, g)$,
where $u$ is a point of $\UUU$
and $g$ is an isomorphism from $X_u$ to $Y_{\bp}$.
An isomorphism $g$ from $X_u$ to $Y_{\bp}$
induces an isomorphism $\bar{g}\colon Q \isom {\sfQ}_{\bp}$.
Conversely,
suppose that an isomorphism 
$\gamma\colon Q \isom {\sfQ}_{\bp}$ is given.
Let $u\in \UUU$ be the point such that $B_{u}=\gamma\inv(\sfB_{\bp})$.
Then $\gamma$ lifts to exactly \emph{two} isomorphisms 
from $X_u$ to $Y_{\bp}$.
Since the variety 
of isomorphisms from $Q $ to ${\sfQ}_{\bp}$
is smooth and irreducible, 
 the fiber of $\varpi_2$ 
over $\bp$ has at most two connected components.
Therefore  $\III$  has at most two connected components,
and hence 
the index $[W: \Gamma]$ of the image
$\Gamma:=\Image(\mon)$ of the monodromy $\mon$ 
in $W:=W(R(X_b))$  is at most $2$.
We assume 
$[W: \Gamma]=2$, 
and derive a contradiction.
\par
Note  that the Weyl group $W$ of type $E_8$ has the  structure $2.G.2$,
where 
\[
2.G=\Ker\; (\det \colon W\to \{\pm 1\}),
\quad 
G.2=W/\{\pm \id\}, 
\]
and $G$ is a simple group. 
See~\cite[page 85]{ATLAS1985}.
By the assumption $[W: \Gamma]=2$,
we have $|\Gamma|=|2.G|=|G.2|$.
Let $\Delta$ be the set of $(-2)$-vectors 
in the root lattice $R(X_b)$. 
For $r\in \Delta$,
let $s_r\in W$ denote the reflection with respect to $r$.
\par
There exists a member $B_q$  of $|\OOO_{Q}(3)|$
that does not pass through $V$ and  has an ordinary node as its only singularity.
By the argument using the local Picard--Lefschetz formula
as in the proof of Proposition~\ref{prop:WE7},
we see that there exists a vanishing cycle
$v\in \Delta$ corresponding to the ordinary double point of $X_q$
such that $\Gamma$ contains the reflection $s_v$.
Since $\det s_v=-1$,
we see that $\Gamma\cap (2.G)$ is of index $2$ in $2.G$.
If $-\id \in \Gamma$, then $-\id\in \Gamma\cap (2.G)$ 
and hence $(\Gamma\cap (2.G))/\{\pm\id\}$ would be a  subgroup of $(2.G)/\{\pm\id\}=G$ with index $2$.
Since $G$ is simple, we have  $-\id \notin \Gamma$.
Hence $\Gamma$ is mapped isomorphically  onto $G.2=W/\{\pm \id\}$,
which acts on $\Delta/\{\pm \id\}$ transitively.
Since $s_v\in \Gamma$ and $g\inv\cdot s_v\cdot g=s_{v^g}=s_{-v^g}$, 
we see that $s_r\in \Gamma$  for any $r\in \Delta$.
Thus we obtain $\Gamma=W$,
which is a contradiction.
%
%
\end{proof}
 %
 %
  \section{Lines in a del Pezzo surface of degree one}\label{sec:lines}
%
We choose a \emph{general}  point $b$ of the parameter space $\UUU$
of the family 
$\XXX\to \UUU$ of bi-anti-canonical models of del Pezzo surfaces of degree $1$ 
treated in Section~\ref{subsec:bianticanonical}.
The purpose of this section is to 
investigate  
the configuration of lines  in the del Pezzo surface  $X_b$.
\par
In Section~\ref{subsec:Bertini}, 
we introduce the notions of 
\emph{Bertini involution} and \emph{tangent plane sections}.
In Section~\ref{subsec:orbitdecomp}, using Proposition~\ref{prop:WE8}, 
we describe the orbit decompositions
of the set of lines in $X_b$ by the monodromy action.
In Section~\ref{subsec:linesinbP2}, we describe the lines in $X_b$ as plane curves on 
the projective plane ${\bfP}^{2}$ obtained by contracting $8$ disjoint lines in $X_b$.
In Section~\ref{subsec:general}, 
we confirm that 
the union of lines in $X_b$ has only ordinary double points as its singularities.
%
\subsection{Bertini involution}\label{subsec:Bertini}
%
We have an orthogonal direct-sum decomposition 
%\[
$\Pic(X_b)=\ZZ \alpha_b \oplus R(X_b)$,
%\Pic(X_b)=\ZZ \alpha_b \oplus R(X_b),
%\]
where $\alpha_b$ is the anti-canonical class of $X_b$.
The orthogonal projection from $\Pic(X_b)$ to $R(X_b)$ induces a bijection 
%
\begin{equation}\label{eq:LX}
L(X_b)\cong  \Delta(R(X_b))
\end{equation}
%
between  the set $L(X_b)$ of lines in $X_b$ and  the set $\Delta(R(X_b))$
of $(-2)$-vectors of the root lattice $R(X_b)$ of type $E_8$.
For a line $l\in L(X_b)$, we denote by 
\[
[l]_R:=[l]-\alpha_b \;\in\; \Delta(R(X_b))
\]
the corresponding $(-2)$-vector. 
Let 
%\[
$\varphi\colon X_b\to Q$
%\] 
be the bi-anti-canonical model of $X_b$,
where $Q\subset \PP^3$ is a   quadric cone with the vertex $V\in Q$,
and let $B_b\subset Q$ be the curve component of  the branch locus of 
$\varphi$.
Then  $B_b$ is a smooth member of $|\OOO_{Q}(3)|$. 
The deck-transformation of  $\varphi$
is called the \emph{Bertini involution},  and is denoted by
\[
\involB\colon X_b\isom X_b.
\]
We call a pair $\{l, l\sprime\}$ of lines in $X_b$ an \emph{$\involB$-pair}
if $l\sprime=\involB(l)$,
and say that $l\sprime=\involB(l)$ is the \emph{$\involB$-partner of $l$}.
It is easy to see that,
for lines $l, l\sprime$ in $X_b$,  the following are equivalent:
(i) the pair $\{l, l\sprime\}$ is an $\involB$-pair, 
(ii)~$\intf{l, l\sprime}=3$, 
(iii)~$[l]+[l\sprime]=2\alpha_b$, and
(iv)~$[l]_R+[l\sprime]_R=0$.
%
%
\begin{definition}\label{def:tps}
A plane section $H\cap Q$ of $Q$,
where $H$ is a linear plane in $\PP^3$,
is called a \emph{tangent plane section for $B_b$}
if $H$ does not pass through the vertex $V$ 
% We need this condition.
%Consider the case where $H$ is a double of a ruling.
and the local intersection number at each intersection point 
of $H$ and $B_b$ is even.
We put
%
\[
S(B_b):=\textrm{the set of 
tangent plane sections  for $B_b$}.
\]
\end{definition}
%
The image $\varphi(l)$ of a line $l\subset X_b$
by $\varphi$ is a  tangent plane section for $B_b$.
Conversely, the pullback by $\varphi$ of a tangent plane section 
is the union of a line and its $\involB$-partner.
Hence we have natural identifications 
%
\begin{equation}\label{eq:barLX}
\barL(X_b):=L(X_b)/\angs{\involB}\;\;\cong\;\;  
\barDelta(R(X_b)):=\Delta(R(X_b))/\{\pm \id\} \;\;\cong\;\; S(B_b).
\end{equation}
%
In particular,
there exist
exactly $120$ tangent plane sections for $B_b$. 
%
%
\begin{remark}\label{rem:thetacharacteristics}
The smooth $(2,3)$-complete intersection  $B_b\subset \PP^3$
is the canonical model of a genus $4$ curve with a vanishing theta constant, and 
the tangent plane sections for $B_b$ are in one-to-one correspondence 
with the odd theta-characteristics of $B_b$.
See Chapter IV and Appendix~B of~\cite{ACGH1985}.
\end{remark}
%
%%%%%%%%%%%%
%
%
\subsection{Orbit decomposition by the monodromy}\label{subsec:orbitdecomp}
%
By Proposition~\ref{prop:WE8},  we can compute 
the monodromy actions of $\pione(\UUU, b)$
on the sets in~\eqref{eq:LX} and~\eqref{eq:barLX}
explicitly.
We describe the orbit decompositions of 
$L(X_b)^{\{k\}}$ and $\barL(X_b)^{\{k\}}$ 
under these monodromy actions.
%
%\par
%\medskip 
\subsubsection{The action on  \texorpdfstring{$L(X_b)^{\{k\}}$}{LXbk}}
The monodromy action on the set $L(X_b)^{\{k\}}$ of $k$-element subsets of $L(X_b)$ 
for small $k$ is as follows.
The numbers of orbits are given as follows:
%
%gap> Read("Task250703c.txt");
%1 1 
%2 4 
%3 12 
%4 62 
%5 378 
%6 3557 
%7 45282 
%
\begin{equation*}%\label{eq:Nks}
\begin{array}{c|ccccccc}
k & 1 & 2 & 3 & 4 & 5 & 6 &7  \\
\hline 
 & 1& 4 & 12 & 62 & 378 &  3557 & 45282\rlap{\;.}
\end{array}
\end{equation*}
%
%gap> Read("Task250610d.txt");
%1 transitive 
%1  done ________________ 
%[ 0, 6720 ] 
%[ 1, 15120 ] 
%[ 2, 6720 ] 
%[ 3, 120 ] 
%2  done ________________ 
%[ [ 0, 0, 0 ], 60480 ] 
%[ [ 0, 0, 1 ], 181440 ] 
%[ [ 0, 0, 2 ], 6720 ] 
%[ [ 0, 1, 1 ], 483840 ] 
%[ [ 0, 1, 2 ], 362880 ] 
%[ [ 0, 2, 2 ], 181440 ] 
%[ [ 0, 2, 3 ], 13440 ] 
%[ [ 1, 1, 1 ], 302400 ] 
%[ [ 1, 1, 2 ], 483840 ] 
%[ [ 1, 1, 3 ], 15120 ] 
%[ [ 1, 2, 2 ], 181440 ] 
%[ [ 2, 2, 2 ], 2240 ] 
%3  done ________________ 
%
%
\begin{itemize}[itemsep=5pt]
%
\item The action on $L(X_b)^{\{1\}}=L(X_b)$  is transitive.
%
\item
The  action on $L(X_b)^{\{2\}}$  has four orbits.
For an orbit $o\subset L(X_b)^{\{2\}}$,
let $m(o)$ denote the intersection number $\intf{l_1, l_2}$ of lines,   
where  $\{l_1, l_2\}\in o$.
Then the four orbits are distinguished  by $m(o)$ as follows:
\[
\begin{array}{c|cccc}
m(o) & 0 & 1 & 2 & 3 \\
\hline 
|o| & 6720 & 15120 & 6720 & 120\mystruth{10pt}
\end{array}.
\]
%
\item
The  action on $L(X_b)^{\{3\}}$  has $12$ orbits.
For an orbit $o\subset L(X_b)^{\{3\}}$,
let $t(o)$ denote 
the non-decreasing sequence of the intersection numbers $\intf{l_i, l_j}$ 
for $1\le i<j\le 3$,
where  $\{l_1, l_2, l_3\}\in o$.
Then the $12$ orbits are described as follows:
%
\begin{equation*}\label{eq:ots}
%\begin{array}{ll}
%t(o) & |o| \\
%\hline 
% {[} 0, 0, 0 {]}& 60480\\
% {[} 0, 0, 1 {]}& 181440\\ 
% {[} 0, 0, 2 {]}& 6720\\
% {[} 0, 1, 1 {]}& 483840\\
% {[} 0, 1, 2 {]}& 362880\\
% {[} 0, 2, 2 {]},& 181440\\
%
%\end{array}
%%
%\qquad
%\begin{array}{ll}
%t(o) & |o| \\
%\hline
% {[} 0, 2, 3 {]}&13440 \\
%{[} 1, 1, 1 {]}& 302400 \\
%{[} 1, 1, 2 {]}& 483840 \\
%{[} 1, 1, 3 {]}&15120 \\
%{[} 1, 2, 2 {]}& 181440 \\
%{[} 2, 2, 2 {]}& 2240 \;\;\llap{.}\\
%%&\\
%%&\\
%\end{array}
\begin{array}{ll}
t(o) & |o| \\
\hline 
 {[} 0, 0, 0 {]}& 60480\\
 {[} 0, 0, 1 {]}& 181440\\ 
 {[} 0, 0, 2 {]}& 6720
 \end{array}
 \qquad
\begin{array}{ll}
t(o) & |o| \\
\hline
 {[} 0, 1, 1 {]}& 483840\\
 {[} 0, 1, 2 {]}& 362880\\
 {[} 0, 2, 2 {]},& 181440
\end{array}
%
\qquad
\begin{array}{ll}
t(o) & |o| \\
\hline
 {[} 0, 2, 3 {]}&13440 \\
{[} 1, 1, 1 {]}& 302400 \\
{[} 1, 1, 2 {]}& 483840 \\
\end{array}
\qquad
\begin{array}{ll}
t(o) & |o| \\
\hline
{[} 1, 1, 3 {]}&15120 \\
{[} 1, 2, 2 {]}& 181440 \\
{[} 2, 2, 2 {]}& 2240 \;\;\llap{.}
\end{array}
\end{equation*}
%
\end{itemize}
%
Let $\LLL\spset{k}\to \UUU$ be the \'etale covering whose fiber over $u\in \UUU$
is $L(X_u)\spset{k}$.
%
\begin{corollary}\label{cor:irred}
{\rm (1)}
The space  $\LLL\spset{2}$ consists of  exactly $4$ irreducible components
$\LLL\spset{2}_0, \dots, \LLL\spset{2}_3$, where 
 $\LLL^{\{2\}}_m\to \UUU$ be the family of 
pairs $\{l_1, l_2\}$ of lines in $X_u$ 
such that $\intf{l_1, l_2}=m$.  %\par
{\rm (2)}
The space $\LLL\spset{3}$ consists of exactly $12$ irreducible components
$\LLL\spset{3}_t$, where 
 $t=[t_1, t_2, t_3]$ runs through the list
%
\begin{equation}\label{eq:tenlist}
\begin{array}{l}
 {[}0,0,0{]},\;\;[0,0,1],\;\;[0,0,2],\;\;[0,1,1],\;\;[0,1,2],\;\; [0,2,2],\;\; [0,2,3], \\
 {[}1,1,1{]},\;\; [1,1,2],\;\; [1,1,3],\;\; [1,2,2],\;\; [2,2,2].
\end{array}
\end{equation}
%
The \'etale covering  $\LLL^{\{3\}}_t\to \UUU$ is the family of 
triples $\{l_1, l_2, l_3\}$ of lines in $X_u$
such that $[\intf{l_1, l_2}, \intf{l_2, l_3}, \intf{l_1, l_3}]$ is equal to  $t$ up to order.
\qed
\end{corollary}
%
\subsubsection{The action on  \texorpdfstring{$\barL(X_b)^{\{k\}}$}{barLXbk}}
The  action on the set  $\barL(X_b)^{\{k\}}$
is identified with 
the action of $\barW$
on $\barDelta^{\{k\}}$ defined in Section~\ref{subsec:Nk}.
Therefore 
the  numbers of orbits are  $N(k)$ and,  for small $k$,
they  are given in
Table~\ref{eq:Nks}.
%
\begin{itemize}[itemsep=5pt]
\item The  action on the set $\barL(X_b)^{\{1\}}=\barL(X_b)$ of $\involB$-pairs is transitive.
%
\item
%
%gap> Read("Task250609d.txt");
%[ 1, 2 ] 3360 20 
%[ 1, 15 ] 3780 11 
%_________________ 
%[ 1, 15, 74 ] 37800 [ 11, 11, 11 ] 20-conn comps 6 
%[ 1, 2, 3 ] 30240 [ 20, 20, 20 ] 20-conn comps 1 
%[ 1, 2, 14 ] 90720 [ 20, 11, 20 ] 20-conn comps 2 
%[ 1, 2, 36 ] 120960 [ 20, 11, 11 ] 20-conn comps 4 
%[ 1, 2, 120 ] 1120 [ 20, 20, 20 ] 20-conn comps 2 
%
The  action decomposes  $\barL(X_b)^{\{2\}}$ into two orbits 
of size $3360 $ and  $3780$.
These two orbits are distinguished by 
Figure~\ref{fig:orbitsbarLX2},
where a line is denoted by a circle 
%
%
%\begin{tikzpicture}[x=1mm, y=1mm]
%\draw[fill=white, draw=black,  thick] (0,0) circle (1);
%\end{tikzpicture},
\includegraphics{smallcircle.pdf},
%
%\newcommand{\iBpair}{\begin{tikzpicture}[x=1mm, y=1mm]
%\draw [line width=2pt] (0,0)--(6,0);
%\draw[fill=white, draw=black,  thick] (0,0) circle (1);
%\draw[fill=white, draw=black,  thick] (6,0) circle (1);
%\end{tikzpicture}}
\newcommand{\iBpair}{\includegraphics{iBpair.pdf}}
%
an $\involB$-pair  is denoted by 
$\iBpair$, 
and the intersection number of distinct two lines is given by 
the number of line-segments connecting the corresponding circles.
%
\item
The action decomposes  $\barL(X_b)^{\{3\}}$ into five orbits.
These  orbits are depicted in
Figure~\ref{fig:orbitsbarLX3}.
%
\end{itemize}
%
%%%%%%%%%%%%%%%%%%%
%
\begin{figure}
\begin{center}
%\begin{tikzpicture}[x=1.1mm, y=1.1mm]
%%
%\begin{scope}[shift={(0,0)}]
%\draw[double, line width=.5pt, double distance=1.5pt]  (0,0)--(15,0);
%\draw[double, line width=.5pt, double distance=1.5pt]  (0,5)--(15,5);
%\draw [line width=2pt] (0,0)--(0,5);
%\draw [line width=2pt] (15,0)--(15,5);
%\draw[fill=white, draw=black,  thick] (0,5) circle (1);
%\draw[fill=white, draw=black,  thick] (0,0) circle (1);
%\draw[fill=white, draw=black,  thick] (15,5) circle (1);
%\draw[fill=white, draw=black,  thick] (15,0) circle (1);
%\draw (8,-4) node {size $3360$};
%\end{scope}
%%
%%
%\begin{scope}[shift={(30,0)}]
%\draw[line width=.5pt]  (0,0)--(15,5);
%\draw[line width=.5pt]  (0,5)--(15,0);
%\draw[line width=.5pt]  (0,0)--(15,0);
%\draw[line width=.5pt]  (0,5)--(15,5);
%\draw [line width=2pt] (0,0)--(0,5);
%\draw [line width=2pt] (15,0)--(15,5);
%\draw[fill=white, draw=black,  thick] (0,5) circle (1);
%\draw[fill=white, draw=black,  thick] (0,0) circle (1);
%\draw[fill=white, draw=black,  thick] (15,5) circle (1);
%\draw[fill=white, draw=black,  thick] (15,0) circle (1);
%\draw (8,-4) node {size $3780$};
%\end{scope}
%%
%\end{tikzpicture}
\includegraphics{orbits2.pdf}
\end{center}
%
\caption{Orbits in $\barL(X_b)^{\{2\}}$}\label{fig:orbitsbarLX2}
\vskip .7cm
\begin{center}
%\begin{tikzpicture}[x=1.1mm, y=1.1mm]
%%
%\pgfmathsetmacro{\s}{30}
%\pgfmathsetmacro{\ss}{\s+\s}
%\pgfmathsetmacro{\sss}{\ss+\s}
%\pgfmathsetmacro{\ssss}{\sss+\s}
%\pgfmathsetmacro{\sssss}{\ssss+\s}
%
%\begin{scope}[shift={(0,0)}]
%%
%\coordinate (A1) at  (0,5);
%\coordinate (A2) at  (0,10);
%\coordinate (B1) at  (4.33,-2.5);
%\coordinate (B2) at  (8.66, -5);
%\coordinate (C1) at  (-4.33, -2.5);
%\coordinate (C2) at  (-8.66, -5);
%%%%% 11:11:11
%\draw[line width=.5pt] (A1)--(B1);
%\draw[line width=.5pt] (A1)--(B2);
%\draw[line width=.5pt] (A1)--(C1);
%\draw[line width=.5pt] (A1)--(C2);
%\draw[line width=.5pt] (A2)--(B1);
%\draw[line width=.5pt] (A2)--(B2);
%\draw[line width=.5pt] (A2)--(C1);
%\draw[line width=.5pt] (A2)--(C2);
%\draw[line width=.5pt] (B1)--(C1);
%\draw[line width=.5pt] (B1)--(C2);
%\draw[line width=.5pt] (B2)--(C1);
%\draw[line width=.5pt] (B2)--(C2);
%%%%%
%\draw[line width=2pt] (A1)--(A2);
%\draw[line width=2pt] (B1)--(B2);
%\draw[line width=2pt] (C1)--(C2);
%\draw[fill=white, draw=black,  thick] (A1) circle (1);
%\draw[fill=white, draw=black,  thick] (A2) circle (1);
%\draw[fill=white, draw=black,  thick] (B1) circle (1);
%\draw[fill=white, draw=black,  thick] (B2) circle (1);
%\draw[fill=white, draw=black,  thick] (C1) circle (1);
%\draw[fill=white, draw=black,  thick] (C2) circle (1);
%\draw (0,-9) node {size $37800$};
%\end{scope}
%%
%\begin{scope}[shift={(\s,0)}]
%%
%\coordinate (A1) at  (0,5);
%\coordinate (A2) at  (0,10);
%\coordinate (B1) at  (4.33,-2.5);
%\coordinate (B2) at  (8.66, -5);
%\coordinate (C1) at  (-4.33, -2.5);
%\coordinate (C2) at  (-8.66, -5);
%%%%% 11:11:20
%\draw[line width=.5pt] (A1)--(B1);
%\draw[line width=.5pt] (A1)--(B2);
%\draw[line width=.5pt] (A1)--(C1);
%\draw[line width=.5pt] (A1)--(C2);
%\draw[line width=.5pt] (A2)--(B1);
%\draw[line width=.5pt] (A2)--(B2);
%\draw[line width=.5pt] (A2)--(C1);
%\draw[line width=.5pt] (A2)--(C2);
%\draw[double, line width=.5pt, double distance=1.5pt]   (B1)--(C1);
%\draw[double, line width=.5pt, double distance=1.5pt]   (B2)--(C2);
%%%%%
%\draw[line width=2pt] (A1)--(A2);
%\draw[line width=2pt] (B1)--(B2);
%\draw[line width=2pt] (C1)--(C2);
%\draw[fill=white, draw=black,  thick] (A1) circle (1);
%\draw[fill=white, draw=black,  thick] (A2) circle (1);
%\draw[fill=white, draw=black,  thick] (B1) circle (1);
%\draw[fill=white, draw=black,  thick] (B2) circle (1);
%\draw[fill=white, draw=black,  thick] (C1) circle (1);
%\draw[fill=white, draw=black,  thick] (C2) circle (1);
%\draw (0,-9) node {size $120960$};
%\end{scope}
%%
%\begin{scope}[shift={(\ss,0)}]
%%
%\coordinate (A1) at  (0,5);
%\coordinate (A2) at  (0,10);
%\coordinate (B1) at  (4.33,-2.5);
%\coordinate (B2) at  (8.66, -5);
%\coordinate (C1) at  (-4.33, -2.5);
%\coordinate (C2) at  (-8.66, -5);
%%%%% 11:20:20
%\draw[line width=.5pt] (A1)--(B1);
%\draw[line width=.5pt] (A1)--(B2);
%\draw[line width=.5pt] (A2)--(B1);
%\draw[line width=.5pt] (A2)--(B2);
%\draw[double, line width=.5pt, double distance=1.5pt]   (A1)--(C1);
%\draw[double, line width=.5pt, double distance=1.5pt]   (A2)--(C2);
%\draw[double, line width=.5pt, double distance=1.5pt]   (B1)--(C1);
%\draw[double, line width=.5pt, double distance=1.5pt]   (B2)--(C2);
%%%%%
%\draw[line width=2pt] (A1)--(A2);
%\draw[line width=2pt] (B1)--(B2);
%\draw[line width=2pt] (C1)--(C2);
%\draw[fill=white, draw=black,  thick] (A1) circle (1);
%\draw[fill=white, draw=black,  thick] (A2) circle (1);
%\draw[fill=white, draw=black,  thick] (B1) circle (1);
%\draw[fill=white, draw=black,  thick] (B2) circle (1);
%\draw[fill=white, draw=black,  thick] (C1) circle (1);
%\draw[fill=white, draw=black,  thick] (C2) circle (1);
%\draw (0,-9) node {size $90720$};
%\end{scope}
%%
%\begin{scope}[shift={(\sss,0)}]
%%
%\coordinate (A1) at  (0,5);
%\coordinate (A2) at  (0,10);
%\coordinate (B1) at  (4.33,-2.5);
%\coordinate (B2) at  (8.66, -5);
%\coordinate (C1) at  (-4.33, -2.5);
%\coordinate (C2) at  (-8.66, -5);
%%%%% 20:20:20 two cycles
%\draw[double, line width=.5pt, double distance=1.5pt]   (A1)--(B1);
%\draw[double, line width=.5pt, double distance=1.5pt]   (A2)--(B2);
%\draw[double, line width=.5pt, double distance=1.5pt]   (A1)--(C1);
%\draw[double, line width=.5pt, double distance=1.5pt]   (A2)--(C2);
%\draw[double, line width=.5pt, double distance=1.5pt]   (B1)--(C1);
%\draw[double, line width=.5pt, double distance=1.5pt]   (B2)--(C2);
%%%%%
%\draw[line width=2pt] (A1)--(A2);
%\draw[line width=2pt] (B1)--(B2);
%\draw[line width=2pt] (C1)--(C2);
%\draw[fill=white, draw=black,  thick] (A1) circle (1);
%\draw[fill=white, draw=black,  thick] (A2) circle (1);
%\draw[fill=white, draw=black,  thick] (B1) circle (1);
%\draw[fill=white, draw=black,  thick] (B2) circle (1);
%\draw[fill=white, draw=black,  thick] (C1) circle (1);
%\draw[fill=white, draw=black,  thick] (C2) circle (1);
%\draw (0,-9) node {size $1120$};
%\end{scope}
%%
%\begin{scope}[shift={(\ssss,0)}]
%%
%\coordinate (A1) at  (0,5);
%\coordinate (A2) at  (0,10);
%\coordinate (B1) at  (4.33,-2.5);
%\coordinate (B2) at  (8.66, -5);
%\coordinate (C1) at  (-4.33, -2.5);
%\coordinate (C2) at  (-8.66, -5);
%%%%% 20:20:20 one cycles
%\draw[double, line width=.5pt, double distance=1.5pt]   (A1)--(B1);
%\draw[double, line width=.5pt, double distance=1.5pt]   (A2)--(B2);
%\draw[double, line width=.5pt, double distance=1.5pt]   (A1)--(C1);
%\draw[double, line width=.5pt, double distance=1.5pt]   (A2)--(C2);
%\draw[double, line width=.5pt, double distance=1.5pt]   (B1)--(C2);
%\draw[double, line width=.5pt, double distance=1.5pt]   (B2)--(C1);
%%%%%
%\draw[line width=2pt] (A1)--(A2);
%\draw[line width=2pt] (B1)--(B2);
%\draw[line width=2pt] (C1)--(C2);
%\draw[fill=white, draw=black,  thick] (A1) circle (1);
%\draw[fill=white, draw=black,  thick] (A2) circle (1);
%\draw[fill=white, draw=black,  thick] (B1) circle (1);
%\draw[fill=white, draw=black,  thick] (B2) circle (1);
%\draw[fill=white, draw=black,  thick] (C1) circle (1);
%\draw[fill=white, draw=black,  thick] (C2) circle (1);
%\draw (0,-9) node {size $30240$};
%\end{scope}
%%
%%
%\end{tikzpicture}
%%
\includegraphics{orbits3.pdf}
\end{center}
%
\caption{Orbits in $\barL(X_b)^{\{3\}}$}\label{fig:orbitsbarLX3}
%
\end{figure}
%
%%%%%%%%%%%%%%%%%%%%%%%%%%%%%%%%%%%%%%
%
\subsection{Lines in the birational model \texorpdfstring{${\bfP}^{2}$}{bP2}}\label{subsec:linesinbP2}
%
We choose  disjoint $8$ lines $l_1, \dots, l_8$ in $X_b$
and consider the contraction 
$\beta\colon X_b\to {\bfP}^{2}$
of these lines.
Let $p_i\in {\bfP}^{2}$ be the point  $\beta(l_i)$ for $i=1, \dots, 8$,
and let $h\in \Pic(X_b)$ be the class of the pullback of a line in ${\bfP}^{2}$.
Since
 \[
 3h=(\alpha_b+l_1+ \dots+l_8),
\]
and we have $\ell+\involB(\ell)=2\alpha_b$, it follows that
%\[
$\intf{h, \ell}+\intf{h, \involB(\ell)}=6$ 
%\]
holds for any line $\ell$.
We investigate the images of the lines $\ell$ by $\beta$.
Calculating the intersection numbers with the exceptional lines $l_1, \dots, l_8$ of $\beta$,
we obtain the following.
(See also~\cite[Theorem~26.2]{Manin1986}.)
In the following, the phrase ``$C$ passes through $p\in {\bfP}^{2}$ \emph{once}"
means that $p$  is a smooth point of $C$. 
%
\begin{itemize}[itemsep=5pt]
%
\item There exist exactly $8$ lines $l_i$ with $h$-degree $0$. 
Their $\involB$-partners are of $h$-degree $6$:
the sextic curve  $\beta(\involB(l_i))\subset {\bfP}^{2}$  has a triple  point at $p_i$, 
and double points at the $7$ points in $\{p_1, \dots, \ p_8\}\setminus\{p_i\}$.
%
\item There exist exactly $28$ lines $l_{ij}$ with $h$-degree $1$,
where $1\le i<j\le 8$.
The line $l_{ij}$ is mapped by $\beta$ to the line  in ${\bfP}^{2}$ 
passing through  $p_i$ and $p_j$.
Their $\involB$-partners are  of   $h$-degree $5$:
the quintic curve  $\beta(\involB(l_{ij}))\subset {\bfP}^{2}$ passes through $p_i$ and $p_j$ once, 
and has  double points at the $6$ points  in 
$\{p_1, \dots, \ p_8\}\setminus\{p_i, p_j\}$.
%
\item There exist exactly $56$ lines $l_{\overline{ijk}}$ with $h$-degree $2$,
where $1\le i<j<k\le 8$.
The line $l_{\overline{ijk}}$  is  mapped by $\beta$ to the conic   in ${\bfP}^{2}$ 
passing through the five points 
in $\{p_1, \dots, p_8\}\setminus \{p_i, p_j, p_k\}$.
Their $\involB$-partners are  of $h$-degree $4$:
the quartic curve  $\beta(\involB(l_{\overline{ijk}}))$
has  double points at 
$p_i, p_j, p_k$ and 
 passes through the $5$ points  in $\{p_1, \dots, p_8\}\setminus \{p_i, p_j, p_k\}$ once.
%
\item There exist exactly $56$ lines $l_{i, j}$
with $h$-degree $3$,
where $1\le i, j\le 8$ and $i\ne j$.
We have $\involB(l_{i, j})=l_{j, i}$.
The cubic curve  $\beta(\involB(l_{i, j}))$
  passes through the $6$ points in 
$\{p_1, \dots, \ p_8\}\setminus\{p_i, p_j\}$ once, 
has a double point at $p_i$, and does not pass through $p_j$.
%
\end{itemize}
%
%See Table~\ref{table:hdegandmults} for the summary of the curves $\beta(\ell)$.
%
%[ [ [ 0, [ [ -1, 1 ], [ 0, 7 ] ] ], 8 ],
% [ [ 1, [ [ 0, 6 ], [ 1, 2 ] ] ], 28 ],
%  [ [ 2, [ [ 0, 3 ], [ 1, 5 ] ] ], 56 ], 
%  [ [ 3, [ [ 0, 1 ], [ 1, 6 ], [ 2, 1 ] ] ], 56 ], 
%  [ [ 4, [ [ 1, 5 ], [ 2, 3 ] ] ], 56 ], 
%  [ [ 5, [ [ 1, 2 ], [ 2, 6 ] ] ], 28 ], 
%  [ [ 6, [ [ 2, 7 ], [ 3, 1 ] ] ], 8 ] ] 
 %
 \begin{table}
 \[
 \begin{array}{ccc}
 \textrm{$h$-degree $\intf{h, \ell}$}  & \textrm{multiplicities $\intf{l_i, \ell}$} & \textrm{number}\\
 \hline
 0 & (-1)^1 0^7 & 8 \mystruth{12pt}\\
1 & 0^6 1^2 & 28 \\
2 &  0^3 1^5 & 56 \\
3 & 0^1 1^6 2^1 & 56 \\
4 & 1^5 2^3 & 56 \\
5 & 1^2 2^6 & 28 \\
6 & 2^7 3^1& 8 
 \end{array}
 \]
 \caption{$240$ curves in \texorpdfstring{${\bfP}^{2}$}{P2}}\label{table:hdegandmults}
 \end{table}
 %
%\begin{remark}
%If we are given the coordinates of the points $p_1, \dots, p_8$
%satisfying the conditions in Definition~\ref{def:PPPn},
%then we can calculate the defining equations of the images  $\beta(\ell)$
% by solving \emph{linear} equations.
%\end{remark}
%
%\begin{remark}
%In~\cite{Degtyarev2012Bertini},
%an explicit description of 
%the Bertini involution $\involB$ on $\bP^2$  is given.
%\end{remark}
 %
\subsection{Union of lines}\label{subsec:general}
%
In this section, we confirm the following result,
which must be well known,  but of which we could not find a reference.
%
\begin{proposition}\label{prop:union}
The union of the $240$ lines  in $X_b$ has only ordinary points 
as its singularities.
\end{proposition}
%
\begin{proof}
Recall that $X_b$ is a \emph{general} member of the family $\XXX\to \UUU$.
By Corollary~\ref{cor:irred},
it is enough to prove the following.
\begin{itemize}
\item[($m$)] For $m=2$ and $m=3$,
there exist a point $u\in \UUU$
and  lines $\ell_1, \ell_2$ in $X_u$ such that
$\intf{\ell_1, \ell_2}=m$, and that 
$\ell_1$ and $\ell_2$ intersect at distinct $m$ points.
\item[(\,$t$\,)]
For each $t=[t_1, t_2, t_3]$ in the second line of~\eqref{eq:tenlist},
there exist a point $u\in \UUU$
and  lines $\ell_1, \ell_2, \ell_3$ in $X_u$ such that $[\intf{\ell_1, \ell_2}, \intf{\ell_2, \ell_3},  \intf{\ell_1, \ell_3}]$ 
is equal to $t$ up to order,
and that  $\ell_1\cap  \ell_2\cap  \ell_3$  is empty.
\end{itemize}
%
We find such a del Pezzo surface $X_u$ by choosing $8$ points 
$p_1, \dots, p_8$ on ${\bfP}^{2}$ satisfying  the conditions in Definition~\ref{def:PPPn}.
Let $Y_{\bp}\to {\bfP}^{2}$ be the blowing-up at these points.
The lines in  $Y_{\bp}$ can be calculated 
by the description given in Section~\ref{subsec:linesinbP2},
and we search for lines satisfying
 the conditions  in  ($m$) and ($t$).
It is enough to find an example over a finite field.
 \par
 We give an example over  $\FF_{19}$.
 We choose the following 
$8$ points:
% newexamples19[1];
%[ [ 0*Z(19), 0*Z(19) ], [ Z(19)^0, 0*Z(19) ], [ 0*Z(19), Z(19)^0 ], [ Z(19)^0, Z(19)^0 ],
% [ Z(19), Z(19)^11 ], [ Z(19)^11, Z(19)^2 ], [ Z(19)^12, Z(19)^11 ],   [ Z(19)^15, Z(19)^4 ] ]
%gap> Z(19)-2*One(GF(19));
%0*Z(19)
%gap> IsZero(Z(19)-2*One(GF(19)));
%true
%gap> [2,2^11] mod 19;
%[ 2, 15 ]
%gap> [2^11,2^2] mod 19;
%[ 15, 4 ]
%gap> [2^12, 2^11] mod 19;
%[ 11, 15 ]
%gap> [2^15, 2^4] mod 19;
%[ 12, 16 ]
 \[
 \begin{aligned}
&p_1=(0,0),
\quad
p_2=(1,0),
\quad
p_3=(0,1),
\quad
p_4=(1,1),\\
\quad
&p_5=(2,15),
\quad
p_6=(15,4),
\quad
p_7=(11,15),
\quad
p_8=(12, 16),
\end{aligned}
 \]
where we use  affine coordinates of ${\bfP}^{2}$.
It is easy to confirm that these points 
satisfy  the conditions in Definition~\ref{def:PPPn}.
A line $\ell$ in $Y_{\bp}$ is denoted as $[d; \mu_1, \dots, \mu_8]$,
where
$d$ is the $h$-degree and $\mu_i$ is the multiplicity $\intf{l_i, \ell}$ of $\beta(\ell)$ at $p_i$.
%Note that, if $\beta(\ell)=\{p_i\}$, then we put $\mu_i=-1$.
We consider the following lines ${\ell}_i$, and calculate the defining equations of $\beta({\ell}_i)$ in ${\bfP}^{2}$:
\[
\begin{array}{lcl}
{\ell}_1 &:=& [0;  -1, 0, 0, 0, 0, 0, 0, 0 ], \\ %[ -1, 0, 0, 0, 0, 0, 0, 0 ] 1 
{\ell}_2 &:=& [3;  2, 1, 1, 1, 1, 1, 1, 0 ], \\ %[ 2, 1, 1, 1, 1, 1, 1, 0 ] 148 
{\ell}_3 &:=& [6;  3, 2, 2, 2, 2, 2, 2, 2 ], \\  %[ 3, 2, 2, 2, 2, 2, 2, 2 ] 240 
\end{array}
\quad
\begin{array}{lcl}
{\ell}_4 &:=& [1; 1, 1, 0, 0, 0, 0, 0, 0 ], \\ %[ 1, 1, 0, 0, 0, 0, 0, 0 ] 36 
{\ell}_5 &:=& [2; 1, 0, 1, 1, 1, 1, 0, 0 ], \\ %[ 1, 0, 1, 1, 1, 1, 0, 0 ] 88 
{\ell}_6 &:=& [3; 2, 0, 1, 1, 1, 1, 1, 1 ], \\ %[ 2, 0, 1, 1, 1, 1, 1, 1 ] 121 
\end{array}
%
\quad
\begin{array}{lcl}
{\ell}_7 &:=& [2; 1, 1, 1, 1, 1, 0, 0, 0 ], \\%[ 1, 1, 1, 1, 1, 0, 0, 0 ] 92 
{\ell}_8 &:=& [6; 2, 3, 2, 2, 2, 2, 2, 2 ], \\% [ 2, 3, 2, 2, 2, 2, 2, 2 ] 239 
{\ell}_9 &:=& [6; 2, 2, 3, 2, 2, 2, 2, 2 ]. \\%[ 2, 2, 3, 2, 2, 2, 2, 2 ] 238
\end{array}
\]
%Then we confirm that 
%\begin{itemize}
%\item the pair $l_1, l_2$ satisfies condition $(m)$ for $m=2$,  %[ 1, 148 ] [ 2 ] 
%\item the pair $l_1, l_3$ satisfies condition $(m)$ for $m=3$, %[ 1, 240 ] [ 3 ] 
%\item the triple $l_1, l_4, l_5$ satisfies condition $(t)$ for $t=[1,1,1]$, %[ 1, 36, 88 ] [ 1, 1, 1 ] 
%\item the triple $l_1, l_4, l_6$ satisfies condition $(t)$ for $t=[1,1,2]$, %[ 1, 36, 121 ] [ 1, 1, 2 ] 
%\item the triple $l_1, l_7, l_3$ satisfies condition $(t)$ for $t=[1,1,3]$, %[ 1, 92, 240 ] [ 1, 1, 3 ] 
%\item the triple $l_1, l_6, l_9$ satisfies condition $(t)$ for $t=[1,2,2]$,  %[ 1, 121, 238 ] [ 1, 2, 2 ] 
%\item the triple $l_1, l_6, l_8$ satisfies condition $(t)$ for $t=[2,2,2]$. %[ 1, 121, 239 ] [ 2, 2, 2 ] 
%\end{itemize}
%
Then we confirm that the pair ${\ell}_1, {\ell}_2$ (resp.~the pair  ${\ell}_1, {\ell}_3$) satisfies condition $(m)$ for $m=2$
 (resp.~$m=3$).
Moreover, 
the triple $\tau=\{{\ell}_i, {\ell}_j, {\ell}_k\}$ satisfies condition $(t)$ for $t=[t_1, t_2, t_3]$,
where
\[
\begin{array}{cc}
\tau & t \\
\hline 
{\ell}_1, {\ell}_4, {\ell}_5 & [1,1,1]\\
{\ell}_1, {\ell}_4, {\ell}_6 & [1,1,2]\\
\end{array}
\quad
\begin{array}{cc}
\tau & t \\
\hline 
{\ell}_1, {\ell}_3, {\ell}_7 & [1,1,3]\\
{\ell}_1, {\ell}_6, {\ell}_9& [1,2,2]\\
\end{array}
\quad
\begin{array}{cc}
\tau & t \\
\hline 
{\ell}_1, {\ell}_6, {\ell}_8 & [2,2,2]\rlap{\;.}\\
 & 
\end{array}
\]
%
%gap> FindPosByMults([-1,0,0,0,0,0,0,0]);
%[ -1, 0, 0, 0, 0, 0, 0, 0 ] 1 
%gap> FindPosByMults([2,1,1,1,1,1,1,0]);
%[ 2, 1, 1, 1, 1, 1, 1, 0 ] 148 
%gap> FindPosByMults([3,2,2,2,2,2,2,2]);
%[ 3, 2, 2, 2, 2, 2, 2, 2 ] 240 
%gap> FindPosByMults([1,1,0,0,0,0,0,0]);
%[ 1, 1, 0, 0, 0, 0, 0, 0 ] 36 
%gap> FindPosByMults([1,0,1,1,1,1,0,0]);
%[ 1, 0, 1, 1, 1, 1, 0, 0 ] 88 
%gap> FindPosByMults([2,0,1,1,1,1,1,1]);
%[ 2, 0, 1, 1, 1, 1, 1, 1 ] 121 
%gap> FindPosByMults([1,1,1,1,1,0,0,0]);
%[ 1, 1, 1, 1, 1, 0, 0, 0 ] 92 
%gap> FindPosByMults([2,3,2,2,2,2,2,2]);
%[ 2, 3, 2, 2, 2, 2, 2, 2 ] 239 
%gap> FindPosByMults([2,2,3,2,2,2,2,2]);
%[ 2, 2, 3, 2, 2, 2, 2, 2 ] 238 
%gap> IntNumbs([1,148]);
%[ 1, 148 ] [ 2 ] 
%gap> IntNumbs([1,240]);
%[ 1, 240 ] [ 3 ] 
%gap> IntNumbs([1,36,88]);
%[ 1, 36, 88 ] [ 1, 1, 1 ] 
%gap> IntNumbs([1,36,121]);
%[ 1, 36, 121 ] [ 1, 1, 2 ] 
%gap> IntNumbs([1,92,240]);
%[ 1, 92, 240 ] [ 1, 1, 3 ] 
%gap> IntNumbs([1,121,238]);
%[ 1, 121, 238 ] [ 1, 2, 2 ] 
%gap> IntNumbs([1,121,239]);
%[ 1, 121, 239 ] [ 2, 2, 2 ] 
%
%
%The defining equations of $\beta({\ell}_i)$ can be found in~\cite{WE8compdata}.
\end{proof}
%
%The following will be used in Section~\ref{sec:family}.
%
\begin{corollary}\label{cor:general1}
{\rm (1)}
Every tangent plane section $H\cap Q$  for $B_b$ intersects $B_b$ at distinct  three  points.
{\rm (2)} 
Suppose that  $H_1\cap Q$ and $H_2\cap Q$ are distinct  tangent plane sections for $B_b$.
Then  $H_1\cap B_b$ and $ H_2\cap B_b$ are disjoint, 
 and  $H_1\cap H_2\cap Q$ consists of distinct two points.
 {\rm (3)} 
 Suppose that  $H_1\cap Q, H_2\cap Q$, and $H_3\cap Q$ are distinct  tangent plane sections for $B_b$.
Then   $H_1\cap H_2\cap H_3 \cap Q$ is empty.
\qed
\end{corollary}
%
%
%\begin{remark}
%In~\cite{DesjardinsWinter2025} and \cite{LuijkWinter2023},
%possible singular points on the union of the  $240$ lines in
%del Pezzo surfaces  of degree $1$ other than ordinary double points  
%(\emph{generalized Eckardt points}) are extensively studied.
%\end{remark}
%
%%%%%%%%%%%%%%%%%
%
%
\section{Del Pezzo surfaces of degree one and \texorpdfstring{$\mtan{3}$}{mtan3}-sextics}\label{sec:mtan3anddelPezzo}
%
In this section,
we relate  $\mtan{3}$-sextics
with del Pezzo surfaces of degree $1$.
In Section~\ref{subsec:planecurveswithmtan},
we prove Propositions~\ref{prop:t3}
and exhibit the parameter space $\TTT$
of $\mtan{3}$-sextics in the frame $(A, \Lambda)$.
In Section~\ref{subsec:pip},
we describe a birational map $\pi_p$ from 
$Q$ to $\PP^2$,
which gives a proof of Propositions~\ref{prop:t3conics},
and induces a birational map between $\UUU$ and $\TTT$.
%
\subsection{Plane curves with \texorpdfstring{$\mtan{m}$}{tm}-singularity}\label{subsec:planecurveswithmtan}
%
We fix a point $A\in \PP^2$
and a line $\Lambda\subset \PP^2$ passing through $A$.
Let $(x, y)$ be affine coordinates 
of $\PP^2$ such that $A=(0,0)$ and $\Lambda=\{y=0\}$.
We consider a plane curve  $C\subset\PP^2$ of degree $d$
defined by 
\[
f(x, y)=\sum_{\mu+\nu\le d} a_{\mu\nu} x^{\mu} y^{\nu}=0.
\]
%
\begin{proposition}\label{prop:tm}
Suppose that $d\ge 2m$.
The plane curve $C=\{f=0\}$ has a $\mtan{m}$-singularity 
at $A$ with the tangent line $\Lambda$ if and only if the following holds:
%
%\begin{enumerate}[label={\rm (\roman*)}]
%\item 
(i) $a_{\mu\nu}=0$ if $\mu+2\nu<2m$, and 
%\item
(ii)  the following equation has distinct $m$ roots:
\[
a_{0, m}\, z^m+a_{2, m-1} \, z^{m-1} + \dots +a_{2m-2, 2}\, z+ a_{2m, 0}=0.
\]
%with variable $z$ has distinct $m$ roots.
%\end{enumerate}
%
\end{proposition}
%
\begin{proof}
We consider the blowing-up 
%\[
%(u, v)\;\; \mapsto\;\; (x, y)=(u, uv)
$(u, v)\mapsto (x, y)=(u, uv)$
%\]
of $\PP^2$ at $A$.
The strict transform of $\Lambda$ is given by $v=0$,
and it intersects with the exceptional divisor $E=\{u=0\}$
at the point $(u, v)=(0,0)$.
Then $C$ has a $\mtan{m}$-singularity 
at $A$ with the tangent line $\Lambda$ if and only if 
the  total transform
%\[
$f(u, uv)=0$
%\]
of $C$ contains $E$ with multiplicity $m$
and the strict transform
\[
\sum_{\mu+\nu\le d} a_{\mu\nu} u^{\mu+\nu -m} v^{\nu}=0
\]
of $C$ has an ordinary $m$-fold point at $(u, v)=(0,0)$
with each local branch intersecting $E$ transversely.
This  condition is  equivalent to conditions (i) and (ii) 
in the statement.
\end{proof}
%
Proposition~\ref{prop:t3} follows from Proposition~\ref{prop:tm} immediately.
%
%
In the following, 
for a point $t\in \TTT$, we denote by $C_t\subset\PP^2$ the corresponding $\mtan{3}$-sextic.
%
\subsection{Birational map \texorpdfstring{$\pi_p$}{pip}}
\label{subsec:pip}
%
Recall that 
$Q\subset \PP^3$ is  a quadric cone %a singular quadric surface of rank $3$
with the vertex $V\in Q$.
Then $Q$ is ruled by lines passing through $V$.
We choose a smooth point $p\in Q\setminus\{V\}$.
Then the projection from $p$ induces a birational map
\[
\pi_p\colon Q\birat \PP^2.
\]
This birational map $\pi_p$ defines a frame $(A, \Lambda)$ in $\PP^2$ as follows.
\par
We describe $\pi_p$ in detail. See also Figure~\ref{fig:birat}.
Let $r_p\subset Q$ be  the line in the ruling passing through $p$, 
and let  $\tilQ\to Q$ be  the composite of the minimal desingularization of $Q$
and the blowing-up at $p$.
Let $E_V$, $E_p$, and $\tilde{r}_p$ be the exceptional curve over $V$,
the exceptional curve over $p$, and the strict transform of $r_p$ in $\tilQ$.
Then $E_V$, $E_p$, $\tilde{r}_p$ are  smooth rational curves on $\tilQ$
with self-intersection number $-2$, $-1$, $-1$, respectively.
We blow down $\tilde{r}_p$ to a point, and then 
we blow down the image $E_V\sprime$ of $E_V$ to a point.
The resulting surface is 
the target  plane $\PP^2$ of $\pi_p$.
The  curves $E_V$ and  $\tilde{r}_p$  are mapped to a point $A\in \PP^2$,
and the curve $E_p$ is mapped to a line $\Lambda$ passing through $A$.
Every member of the ruling of $Q$ other than $r_p$ is mapped to a line  passing through $A$.
A plane section $H\cap Q$ not containing $V$ and $p$ is mapped to a smooth conic 
that is  passing through $A$ and is tangent to $\Lambda$ at $A$.
\par
The inverse of $\pi_p$ is given as follows.
Let $(\PP^2)^{\mathord{\sim}}\to \PP^2$ be the blowing up at $A$,
and let $\tilQ\to (\PP^2)^{\mathord{\sim}}$  be the blowing up at the intersection point $A\sprime$
of the exceptional curve $E_A$ over $A$ and the strict transform of $\Lambda$.
Then the strict transform $\widetilde{E}_A$ of $E_A$ (resp.~$\widetilde{\Lambda}$ of $\Lambda$) in $\tilQ$ 
is a smooth rational curve of self-intersection number $-2$ (resp.~$-1$).
We contract these two curves, and 
let $\tilQ\to Q$ denote  the contraction map.
Then $\widetilde{E}_A$  is contracted to the singular point $V$ of $Q$,  
and $\widetilde{\Lambda}$ is contracted to the  center $p$
of the projection.
The exceptional curve $E_{A\sprime}\subset \tilQ$ over $A\sprime$ is 
mapped to the line $r_p$ in the ruling.  
%
\begin{figure}
\begin{center}
%\begin{tikzpicture}[x=.4mm, y=.4mm]
%%
%\begin{scope}[scale=.9]
%\draw[thick, rotate=-45] (0,0) ellipse [x radius=30, y radius=15];
%\draw[thick] (-70, -70)--(-10.6, -10.6);
%\draw[thick] (-70, -70)--(-22.8, 18.9);
%\draw[thick] (-70, -70)--(18.9, -22.8);
%\filldraw[black] (-20,-20) circle[radius=1.5];
%\filldraw[black] (-69, -69) circle[radius=1.5];
%\draw (-10.6, -9.5) node[right]{$r_p$};
%\draw (-69.5, -65) node[above]{$V$};
%\draw (-20,-22) node[below]{$p$};
%\end{scope}
%%
%%
%%
%\begin{scope}[shift={(95, -20)}, scale=1.2]
%\draw[->, very thick] (70,25)--(85, 10);
%\filldraw[black] (90, -10) circle[radius=1.5];
%\draw[thick] (60, -10)--(120, -10);
%\draw (90, -10) node[below]{$A$};
%\draw (120, -10) node[right]{$\Lambda$};
%\end{scope}
%%
%\begin{scope}[shift={(30,-15)}, scale=.9]
%\draw[<-, very thick] (0,25)--(15, 40);
%\draw[thick] (50, 100)--(120, 100);
%\draw[thick] (50, 60)--(75, 120);
%\draw[thick] (120, 60)--(95, 120);
%\draw(120, 100) node[right]{$\tilde{r}_p=E_{A\sprime}$};
%\draw(50, 60) node[below]{$E_V=\widetilde{E}_A$};
%\draw(120, 60) node[below]{$E_p=\widetilde{\Lambda}$};
%\draw(120, 125) node[above]{\phantom{aaa}};
%\end{scope}
%%
%\end{tikzpicture}
\includegraphics{birat.pdf}
\end{center}
\caption{Birational map $\pi_p$}\label{fig:birat}
\end{figure}
%
\subsection{Birational map \texorpdfstring{$\Pi_p$}{Pip}}
\label{subsec:Pip}
%
Recall that  $\UUU\subset |\OOO_{Q}(3)|$ is the parameter space of
the family $\XXX\to \UUU$ of bi-anti-canonical models of del Pezzo surfaces of degree $1$, and that 
 $\TTT$ is the parameter space of the family of $\mtan{3}$-sextics $C_t\subset \PP^2$ 
in the frame $(A, \Lambda)$.
Note that both of $\TTT$ and $\UUU$ are of dimension $15$.
\par
We have chosen a point $p\in Q\setminus\{V\}$.
Suppose that $u$ is a general point of $\UUU$.
In particular, the curve $B_u$ does not contain $p$ and the ruling line $r_p$ intersects $B_u$ 
at three distinct points.
Then the image of  $B_u$ by $\pi_p$ is a $\mtan{3}$-sextic $C_{t(u)}$ in the frame $(A, \Lambda)$,
where $t(u)\in \TTT$.
Conversely, if $t$ is a general point of $\TTT$,
then the image of the $\mtan{3}$-sextic $C_t$ by the inverse of $\pi_p$ is a member $B_{u(t)}$ of $\UUU$.
Therefore
the birational map
$\pi_p$ between $Q$ and $\PP^2$ induces a birational map %between $\UUU$ and $\TTT$.
%
\begin{equation}\label{eq:biratUUUTTT}
\Pi_p\colon \UUU\birat\TTT.
\end{equation}
%
%This birational map plays an important role in 
%the proof of  Theorem~\ref{thm:main2} in 
%Section~\ref{sec:embtop}.
%
\begin{proof}[Proof of Proposition~\ref{prop:t3conics}]
Recall that
$G(C_t)$ denotes the set of special tangent conics 
of  the $\mtan{3}$-sextic $C_t$ for $t\in \TTT$, and that
 $S(B_u)$ denotes the set of tangent plane sections for $B_u$ for $u\in \UUU$.
Suppose that $u\in \UUU$ is  general.
Then the  birational map
$\pi_p$ induces a bijection
%\[
%S(B_u)\;\cong\;G(C_{t(u)}).
$S(B_u)\cong G(C_{t(u)})$.
%\]
Moreover, 
the properties of the tangent plane sections for $B_u$ 
given in Corollary~\ref{cor:general1} carry over to properties
of the special tangent conics of $C_t$,
ensuring that they are in a general position.
\end{proof}
 %
 \section{Embedding topology}\label{sec:embtop}
 %
 In this section, we prove Theorem~\ref{thm:main2}.
 The following observation is due to Artal~Bartolo.
 %
 \begin{lemma}\label{lem:AB}
 Any self-homeomorphism of $\PP^2$ preserves the orientation.
 \end{lemma}
 %
  \begin{proof}
  Note that the intersection form on $H_2(\PP^2)$ 
  depends on the choice of an orientation.
  For an orientation $\xi$ of $\PP^2$, let $\intf{\phantom{i,i}}_{\xi}$ denote 
  the corresponding  intersection form on $H_2(\PP^2)$.
  Then we have $\intf{\phantom{i,i}}_{-\xi}=-\intf{\phantom{i,i}}_{\xi}$.
  Let $\eta$ be the orientation coming from the complex structure of $\PP^2$,
which satisfies $\intf{\gamma, \gamma}_{\eta}\ge 0$
  for any  $\gamma\in H_2(\PP^2)$.
Suppose that a self-homeomorphism $f$ of $\PP^2$ satisfies $f_*\eta = -\eta$.
Then we have 
  \[
  \intf{\gamma,\gamma}_{\eta}=  \intf{f_* \gamma,f_* \gamma}_{f_* \eta}=- \intf{f_* \gamma,f_* \gamma}_{\eta}, 
  \]
which is a contradiction.
  \end{proof}
  %
  Now we prove our main result.
 %
 \begin{proof}[Proof of Theorem~\ref{thm:main2}]
 Let $\UUU$ and $\TTT$ be as in Section~\ref{subsec:Pip}.
 We choose Zariski open subsets $\TTT^0\subset \TTT$ and $\UUU^0\subset \UUU$
such that $\TTT^0$ and $\UUU^0$ are isomorphic via the birational map
$\Pi_p\colon \UUU\birat \TTT$,  
and such that
the special tangent conics of $C_t$ are 
in general position for any $t\in \TTT_0$.
We fix a  point $o\in \TTT^0$, 
and let $b\in \UUU^0$ be the  point corresponding to $o$ via $\Pi_p$. 
Consider the family  $\GGG^0\to \TTT^0$
whose fiber over $t\in  \TTT^0$ is the set  $G(C_t)$ of special tangent conics. 
Then $\GGG^0\to \TTT^0$ is 
isomorphic to the pullback 
of the family $\SSS \to \UUU$
of the sets $S (B_u)$ of tangent plane sections 
via the morphism
%
\begin{equation}\label{eq:TTT0UUU}
\TTT^0\;\;\cong\;\; \UUU^0\;\;\inj\;\; \UUU.
\end{equation}
%
In particular,
the monodromy action 
of $\pione(\TTT^0, o)$ on $G(C_o)$ is induced 
by
the monodromy action 
of $\pione(\UUU, b)$ on $S (B_b)$
via the bijection $S (B_b)\cong G(C_o)$ given by $\pi_p$ 
and the surjective homomorphism
\[
\pione(\TTT^0, o)\surj \pione(\UUU, b)
\]
induced by~\eqref{eq:TTT0UUU}.
Recall from Proposition~\ref{prop:WE8} that,
under the identification of $S (B_b)$ with  $\barDelta(R(X_b))$ in~\eqref{eq:barLX}, 
 the monodromy action 
of $\pione(\UUU, b)$ on $S (B_b)$
factors through 
the surjective homomorphism~\eqref{eq:mon}
to the Weyl group $W(R(X_b))$. % of type $E_8$.
\par
We  consider 
the family $\GGG\sp{0\{k\}} \to \TTT^0$
whose fiber over $t$ is the set $G(C_t)\spset{k}$.
By the above discussion on the monodromy and the definition of $N(k)$, 
the space  $\GGG\sp{0\{k\}}$  has exactly $N(k)$ connected components.
If two points $s, s\sprime$
of $G(C_o)\spset{k}$ are in the same connected component 
of $\GGG\sp{0\{k\}}$,
then we can deform $D_{o,s}$ to   $D_{o, s\sprime}$ in $\PP^2$
without changing the embedding topology
along a path in $\GGG\sp{0\{k\}}$ connecting $s$ and $s\sprime$.
\par
To complete the proof of Theorem~\ref{thm:main2}, it is enough to show that, 
if $D_{o,s}$ and $D_{o, s\sprime}$ have the same embedding topology,
then $s$ and $s\sprime$ belong to the same connected component 
of $\GGG\sp{0\{k\}}$.
For a special tangent conic $\Gamma\in G(C_o)$ of $C_o$,
let $\delta_{\Gamma}\in \barDelta(R(X_b))$ 
denote the pair $\{[l_{\Gamma}]_R, -[l_{\Gamma}]_R\}\subset R(X_b)$,
where $\{l_{\Gamma}, \involB(l_{\Gamma})\}$ is the $\involB$-pair  obtained from 
the tangent plane section 
%$H_{\Gamma}\cap Q$ 
for $B_b$
corresponding to $\Gamma$ via $\pi_p$.
We then put
\[
\delta(s):=\set{\delta_{\Gamma}}{\Gamma\in s}\;\;\in\;\; \barDelta(R(X_b))\spset{k}.
\]
To show that $s$ and $s\sprime$ are in the same connected component of $\GGG\sp{0\{k\}}$,
it is enough to 
find an isometry $g\in W(R(X_b))$ of the lattice $R(X_b)$
such that
the self-bijection of $\barDelta(R(X_b))\spset{k}$ induced by $g$ 
maps $\delta(s)$ to $\delta(s\sprime)$.
\par
Suppose that  $D_{o,s}$ and $D_{o, s\sprime}$ have the same embedding topology.
We have a homeomorphism
\[
\Psi\;\;\colon\;\; (\PP^2, D_{o, s})\;\isom\; (\PP^2, D_{o, s\sprime}).
\]
Then $\Psi$ induces a self-homeomorphism $\Psi_M$ of the complement 
\[
M_o:=\PP^2-(C_o+\Lambda).
\]
Since $H_1(M_o)\cong \ZZ$,
there exists a unique double covering 
%\[
$\varphi_M\colon Z_o\to M_o$ 
%\]
of $M_o$
by a connected surface $Z_o$, and  $\Psi_M$ 
 lifts to a self-homeomorphism $\Psi_Z$ of $Z_o$.
 Note that the lift $\Psi_Z$ is unique up to the deck-transformation of $Z_o$ over $M_o$.
Since $\Psi_M$ is the restriction of a self-homeomorphism of $\PP^2$,
Lemma~\ref{lem:AB} implies that   $\Psi_M$
preserves the orientation of $M_o$, 
and hence  $\Psi_Z$
is an orientation-preserving homeomorphism of $Z_o$.
Consequently, the automorphism 
%\[
%g_Z(\Psi)\colon H_2(Z_o)\isom H_2(Z_o)
$g_Z(\Psi)\colon H_2(Z_o)\isom H_2(Z_o)$
%\]
of the $\ZZ$-module $H_2(Z_o)$
induced by $\Psi_Z$  
preserves the intersection form $\intf{\phantom{i,i}}_Z$
given  by the complex structure of $Z_o$.
We put
\[
\Ker\, \intf{\phantom{i,i}}_Z:=\set{x\in H_2(Z_o)}{\intf{x, y}_Z=0\;\textrm{for any}\; y\in H_2(Z_o)}.
\]
Then $g_Z(\Psi)$ gives rise to  
an isometry  
of the lattice 
\[
\barH_2(Z_o):=H_2(Z_o)/\Ker\, \intf{\phantom{i,i}}_Z.
\]
Note that $\pi_p$ induces an isomorphism from $M_o$ to 
\[
Q_b^0:=Q-(B_b+r_p).
\]
Let $2\tilB_b$ and $(r_p)\sp{\sim}$ be the pullback of $B_b$ and $r_p$ by the double covering
$X_b\to Q$. 
Then the double covering $Z_o$ of $M_o$ can be identified with 
\[
X_b^0:=X_b-(\tilB_b+(r_p)\sp{\sim}).
\]
Note that this identification 
 is unique up to the deck-transformation of $X_b^0$ over~$Q_b^0$.
In the following,  
we regard $Z_o$ as a Zariski open subset of $X_b$.
We also  identify $H_2(X_b)$ with $H^2(X_b)=\Pic(X_b)$ by the Poincar\'e duality.
Since the homology classes of $\tilB_b$ and  $(r_p)\sp{\sim}$  in $H_2(X_b)$ are $3\alpha_b$ and $\alpha_b$,
respectively, and both of $\tilB_b$ and  $(r_p)\sp{\sim}$ are irreducible, 
the Poincar\'e-Lefschetz duality implies  that the inclusion $Z_o\inj X_b$ yields a surjective homomorphism
%
\begin{equation}\label{eq:iotahom}
H_2(Z_o)\surj (\alpha_b)\sperp \subset H_2(X_b)
\end{equation}
%
that preserves the intersection form,
where $(\alpha_b)\sperp$ is the orthogonal complement of $\alpha_b$. 
By the identification $H_2(X_b)=\Pic(X_b)$,
we have $(\alpha_b)\sperp=R(X_b)$, and 
hence  the intersection form on $(\alpha_b)\sperp$ is non-degenerate.
Therefore 
the kernel of~\eqref{eq:iotahom} is equal to $\Ker \intf{\phantom{i,i}}_Z$.
In particular, the lattice $\barH_2(Z_o)$ is  isomorphic to the lattice $R(X_b)$.
Consequently,  the homeomorphism $\Psi$ defines an isometry
\[
g_X(\Psi)\;\in\; W(R(X_b)).
\]
Note that $g_X(\Psi)$ is uniquely determined by $\Psi$  up to $\pm \id$.
%
%
\par
We will show that 
the self-bijection of $\barDelta(R(X_b))\spset{k}$ induced by 
 $g_X(\Psi)$ 
sends $\delta(s)=\set{\delta_{\Gamma}}{\Gamma\in s}$ to $\delta(s\sprime)$.
Since $\Psi$ maps the elements of $s$ to the elements of $s\sprime$ bijectively,
it suffices to show that,
if $\Psi$ maps a special tangent conic $\Gamma$ to a special tangent conic $\Gamma\sprime$,
then $g_X(\Psi)$ maps $\delta_{\Gamma}$ to $\delta_{\Gamma\sprime}$.
\par
Suppose that $\Psi(\Gamma)=\Gamma\sprime$.
The curve $\varphi_M\inv(\Gamma\cap M_o)$
in $Z_o$ has two connected components $\Gamma^+$ and $\Gamma^-$.
The closure of $\Gamma^+$ in $X_b$ is a line $l_{\Gamma}$,
and the closure of $\Gamma^-$ is its $\involB$-partner.
(Recall that we have $Z_o\subset X_b$.)
These two components, viewed as locally finite topological cycles, 
give rise to
linear forms
\[
\gamma^+\colon H_2(Z_o)\to \ZZ,
\quad 
\gamma^-\colon H_2(Z_o)\to \ZZ, 
\]
by the intersection paring.
By definition, 
each of  these linear forms  factors through the quotient homomorphism 
\[
H_2(Z_o)\surj (\alpha_b)\sperp= R(X_b)
\] 
given by $Z_o\inj X_b$,
and the induced  linear form $(\alpha_b)\sperp\to \ZZ$  coincides with  the intersection pairing with $[l_{\Gamma}]_R\in R(X_b)$
and  with $[\involB(l_{\Gamma})]_R=-[l_{\Gamma}]_R$, respectively.
The connected components of 
the curve $\varphi_M\inv(\Gamma\sprime \cap M_o)$ are
$\Psi_Z(\Gamma^{+})$ and $\Psi_Z(\Gamma^{-})$.
The linear forms
on $H_2(Z_o)$ given by these locally finite topological cycles
are equal to
$\gamma^{+}\circ g_Z(\Psi)\inv$ and $\gamma^{-}\circ g_Z(\Psi)\inv $, respectively.
Each of these linear forms yields
a linear form on $(\alpha_b)\sperp= R(X_b)$,
which is the intersection pairing  with $ [l_{\Gamma\sprime}]_R$ and with $-[l_{\Gamma\sprime}]_R$,
respectively, where $l_{\Gamma\sprime}$ is the closure of $\Psi_Z(\Gamma^{+})$ in $X_b$.
Hence we obtain $g_X(\Psi)( [l_{\Gamma}]_R)= [l_{\Gamma\sprime}]_R$.
\end{proof}
%
\begin{remark}\label{rem:K3}
The double covering $W$ of $X_b$ branching along a general member of $|2\alpha_b|$
is a $K3$ surface.
In~\cite{Roulleau2022}, it is proved that the automorphism group of $W$ is isomorphic to $ \ZZ/2\ZZ\times\ZZ/2\ZZ$, 
that $W$ contains exactly  $240$ smooth rational curves, and that 
 $W$ has a structure of the double plane $W\to\PP^2$
whose branch curve is a smooth sextic possessing  $120$  conics that are $6$-tangent.
This $K3$ surface had been discovered in~\cite{Kondo1989}.
\end{remark}



\section*{Acknowledgment}
The author expresses his sincere gratitude to
Professor Enrique Artal~Bartolo for providing Lemma~\ref{lem:AB} and for many valuable discussions.
He also thanks
Professor Shinzo Bannai, Professor Akira Ohbuchi, and Professor Masahiko Yoshinaga
for many comments and discussions, and is grateful to the anonymous referee for valuable comments.
%2026/04/01 "grateful to the anonymous referee"
\par
The author  was partially supported by JSPS KAKENHI Grant Number~20H00112, 23K20209, and 23H00081.
\par
Conflict of Interest: The authors declare that they have no conflict of interest.
\par
During the preparation of this work, the author used ChatGPT (OpenAI, GPT-4.0) in order to improve the clarity and grammar of the English text. 
After using this tool, the author reviewed and edited the content as needed and takes full responsibility for the content of the publication.


\begin{thebibliography}{10}
\bibitem{ACGH1985}
Arbarello E, Cornalba M, Griffiths PA, Harris J.
Geometry of algebraic curves. Vol. I. Grundlehren der mathematischen Wissenschaften, vol. 267.
New York, NY, USA: Springer-Verlag, 1985.

\bibitem{Artal1994}
Artal Bartolo E.
Sur les couples de Zariski.
Journal of Algebraic Geometry 1994; 3 (2): 223-247.

\bibitem{TheSurvey}
Artal Bartolo E, Cogolludo JI, Tokunaga H.
A survey on Zariski pairs.
In: Algebraic geometry in East Asia--Hanoi 2005.
Advanced Studies in Pure Mathematics, vol. 50.
Tokyo, Japan: Mathematical Society of Japan, 2008, pp. 1-100.
\url{https://doi.org/10.2969/aspm/05010001}

\bibitem{Bannai2016}
Bannai S.
A note on splitting curves of plane quartics and multi-sections of rational elliptic surfaces.
Topology and its Applications 2016; 202: 428-439.
\url{https://doi.org/10.1016/j.topol.2016.02.005}

\bibitem{Bannai2020}
Bannai S, Ohno M.
Two-graphs and the embedded topology of smooth quartics and its bitangent lines.
Canadian Mathematical Bulletin 2020; 63 (4): 802-812.
\url{https://doi.org/10.4153/s0008439520000053}

\bibitem{Bannai2019}
Bannai S, Tokunaga H, Yamamoto M.
A note on the topology of arrangements for a smooth plane quartic and its bitangent lines.
Hiroshima Mathematical Journal 2019; 49 (2): 289-302.
\url{https://doi.org/10.32917/hmj/1564106549}

\bibitem{ATLAS1985}
Conway JH, Curtis RT, Norton SP, Parker RA, Wilson RA.
ATLAS of finite groups.
Eynsham, UK: Oxford University Press, 1985.

\bibitem{Demazure1980}
Demazure M.
Surfaces de del Pezzo, II, III, IV, V.
In: Séminaire sur les Singularités des Surfaces.
Lecture Notes in Mathematics, vol. 777.
Berlin, Germany: Springer, 1980, pp. 23-69.

\bibitem{Harris1979}
Harris J.
Galois groups of enumerative problems.
Duke Mathematical Journal 1979; 46 (4): 685-724.
\url{https://doi.org/10.1215/S0012-7094-79-04635-0}

\bibitem{Kondo1989}
Kondo S.
Algebraic $K3$ surfaces with finite automorphism groups.
Nagoya Mathematical Journal 1989; 116: 1-15.
\url{https://doi.org/10.1017/S0027763000001653}

\bibitem{Lamotke1981}
Lamotke K.
The topology of complex projective varieties after S. Lefschetz.
Topology 1981; 20 (1): 15-51.
\url{https://doi.org/10.1016/0040-9383(81)90013-6}

\bibitem{Manin1986}
Manin YI.
Cubic forms: Algebra, geometry, arithmetic.
North-Holland Mathematical Library, vol. 4.
Amsterdam, Netherlands: North-Holland Publishing Co., 1986.

\bibitem{Roulleau2022}
Roulleau X.
An atlas of $K3$ surfaces with finite automorphism group.
Épijournal de Géométrie Algébrique 2022; 6: Article 19, 95 pp.
\url{https://doi.org/10.46298/epiga.2022.6286}

\bibitem{Shimada2003}
Shimada I.
Equisingular families of plane curves with many connected components.
Vietnam Journal of Mathematics 2003; 31 (2): 193-205.

\bibitem{Shimada2010}
Shimada I.
Topology of curves on a surface and lattice-theoretic invariants of coverings of the surface.
In: Algebraic geometry in East Asia--Seoul 2008.
Advanced Studies in Pure Mathematics, vol. 60.
Tokyo, Japan: Mathematical Society of Japan, 2010, pp. 361-382.
\url{https://doi.org/10.2969/aspm/06010361}

\bibitem{Shimada2022}
Shimada I.
Zariski multiples associated with quartic curves.
Journal of Singularities 2022; 24: 169-189.
\url{https://doi.org/10.5427/jsing.2022.24g}

\bibitem{Shimada2026}
Shimada I.
Zariski pairs on cubic surfaces.
Journal of Knot Theory and its Ramifications 2026; 35: Article 2550001.
\url{https://doi.org/10.1142/S021821652550083X}

\bibitem{ZariskiA}
Zariski O.
A theorem on the Poincaré group of an algebraic hypersurface.
Annals of Mathematics 1937; 38 (1): 131-141.
\url{https://doi.org/10.2307/1968515}

\bibitem{ZariskiB}
Zariski O.
The topological discriminant group of a Riemann surface of genus $p$.
American Journal of Mathematics 1937; 59 (2): 335-358.
\url{https://doi.org/10.2307/2371416}

%
%\bibitem{ACGH1985}
%Arbarello E, Cornalba M, Griffiths PA, Harris J.
%Geometry of algebraic curves. {V}ol. {I}, volume 267 of
%  Grundlehren der mathematischen Wissenschaften.
%  New York, NY, USA: Springer-Verlag,1985.
%% \url{https://doi.org/10.1007/978-1-4757-5323-3}.
%
%\bibitem{Artal1994}
%Artal Bartolo E.
% Sur les couples de {Z}ariski.
%{Journal of Algebraic Geometry}, 3(2):223--247, 1994.
%% no doi
%
%%\bibitem{TheSurvey}
%%E.~Artal~Bartolo, J.~I. Cogolludo,  Tokunaga H.
%% A survey on {Z}ariski pairs.
%% In{Algebraic geometry in {E}ast {A}sia---{H}anoi 2005},
%%  volume~50 of{Adv. Stud. Pure Math.}, pages 1--100. Math. Soc. Japan,
%%  Tokyo, 2008.
%% \url{https://doi.org/10.2969/aspm/05010001}
%\bibitem{TheSurvey}
%Artal Bartolo E, Cogolludo JI, Tokunaga H. 
%A survey on Zariski pairs. 
%  In: Algebraic geometry in East Asia--Hanoi 2005. 
% Advanced Studies in Pure Mathematics, vol. 50. Tokyo, Japan: Mathematical Society of Japan, 2008, pp. 1-100. 
% \url{https://doi.org/10.2969/aspm/05010001}
%
%\bibitem{Bannai2016}
%Bannai S.
% A note on splitting curves of plane quartics and multi-sections of
%  rational elliptic surfaces.
%{Topology and its Applications}, 202:428--439, 2016.
% \url{https://doi.org/10.1016/j.topol.2016.02.005}
%
%\bibitem{Bannai2020}
%Bannai S, Ohno M.
% Two-graphs and the embedded topology of smooth quartics and its
%  bitangent lines.
%{Canadian Mathematical Bulletin. Bulletin Canadien de
%  Math\'ematiques}, 63(4):802--812, 2020.
% \url{https://doi.org/10.4153/s0008439520000053}
%
%\bibitem{Bannai2019}
%Bannai S, Tokunaga H, Yamamoto M.
% A note on the topology of arrangements for a smooth plane quartic and
%  its bitangent lines.
%{Hiroshima Mathematical Journal}, 49(2):289--302, 2019.
% \url{https://doi.org/10.32917/hmj/1564106549}
%
%\bibitem{ATLAS1985}
%Conway JH, Curtis RT, Norton SP, Parker RA, Wilson RA.
%%{{$\mathbb{ATLAS}$} of finite groups}.
%{{ATLAS} of finite groups}.
% Eynsham, UK: Oxford University Press, 1985.
%% Oxford University Press, Eynsham, 1985.
%% Maximal subgroups and ordinary characters for simple groups, With
%%  computational assistance from J. G. Thackray.
%
%\bibitem{Demazure1980}
%Demazure M. 
%Surfaces de del Pezzo, II, III, IV, V. In: S\'eminaire sur les Singularit\'es des Surfaces. Lecture Notes in Mathematics, vol. 777. Berlin, Germany: Springer, 1980, pp. 23-69.
%%
%%M.~Demazure.
%% {S}urfaces de del {P}ezzo, {II}, {III}, {IV}, {V}.
%% In{S\'eminaire sur les {S}ingularit\'es des {S}urfaces}, volume
%%  777 of{Lecture Notes in Mathematics}, pages 23 -- 69. Springer, Berlin,
%%  1980.
%% Held at the Centre de Math\'ematiques de l'\'Ecole Polytechnique,
%%  Palaiseau, 1976--1977.
%
%%\bibitem{GAP}
%%The GAP~Group.
%%{{GAP} -- {G}roups, {A}lgorithms, and {P}rogramming, Version
%%  4.13.0}, 2024.
%% \url{https://www.gap-system.org}.
%
%\bibitem{Harris1979}
%Harris J.
% Galois groups of enumerative problems.
%{Duke Mathematical Journal}, 46(4):685--724, 1979.
% \url{https://doi.org/10.1215/S0012-7094-79-04635-0}
%
%\bibitem{Kondo1989}
%Kondo S.
% Algebraic {$K3$} surfaces with finite automorphism groups.
%{Nagoya Mathematical Journal}, 116:1--15, 1989.
% \url{https://doi.org/10.1017/S0027763000001653}
%
%\bibitem{Lamotke1981}
%Lamotke K.
% The topology of complex projective varieties after {S}. {L}efschetz.
%{Topology. An International Journal of Mathematics},
%  20(1):15--51, 1981.
% \url{https://doi.org/10.1016/0040-9383(81)90013-6}
%
%\bibitem{Manin1986}
%%Y.~I. Manin.
%%{Cubic forms. Algebra, geometry, arithmetic}, volume~4 of {\em
%%  North-Holland Mathematical Library}.
%% North-Holland Publishing Co., Amsterdam, second edition, 1986.
%% Translated from the Russian by M. Hazewinkel.
%Manin YI. Cubic forms: Algebra, geometry, arithmetic. 2nd ed. North-Holland Mathematical Library, vol. 4. Amsterdam, Netherlands: North-Holland Publishing Co., 1986. 
%%Translated from the Russian by M. Hazewinkel.
%
%\bibitem{Roulleau2022}
%%Roulleau X.
%% An atlas of {K}3 surfaces with finite automorphism group.
%%{\'Epijournal de G\'eo\'etrie Alg\'ebrique. EPIGA}, 6:Art. 19,
%%  95, 2022.
%%  \url{https://doi.org/10.46298/epiga.2022.6286}
%Roulleau X. An atlas of $K3$ surfaces with finite automorphism group. \'Epijournal de G\'eo\'etrie Alg\'ebrique 2022; 6: Article 19, 95 pp. 
%  \url{https://doi.org/10.46298/epiga.2022.6286}
%
%\bibitem{Shimada2003}
%Shimada I.
% Equisingular families of plane curves with many connected components.
%{Vietnam Journal of Mathematics}, 31(2):193--205, 2003.
%% no doi
%
%\bibitem{Shimada2010}
%%Shimada I.
%% Topology of curves on a surface and lattice-theoretic invariants of
%%  coverings of the surface.
%% In{Algebraic geometry in {E}ast {A}sia---{S}eoul 2008},
%%  volume~60 of{Adv. Stud. Pure Math.}, pages 361--382. Math. Soc. Japan,
%%  Tokyo, 2010.
%% \url{https://doi.org/10.2969/aspm/06010361}
% Shimada I. Topology of curves on a surface and lattice-theoretic invariants of coverings of the surface. In: Algebraic geometry in East Asia--Seoul 2008. Advanced Studies in Pure Mathematics, vol. 60. Tokyo, Japan: Mathematical Society of Japan, 2010, pp. 361-382.
% \url{https://doi.org/10.2969/aspm/06010361}
% 
%
%\bibitem{Shimada2022}
%Shimada I.
% Zariski multiples associated with quartic curves.
%{Journal of Singularities}, 24:169--189, 2022.
% \url{https://doi.org/10.5427/jsing.2022.24g}
%
%\bibitem{Shimada2026}
%Shimada I.
% Zariski pairs on cubic surfaces.
%{Journal of Knot Theory and its Ramifications}, 35:Paper No. 2550001, 2026.
% \url{https://doi.org/10.1142/S021821652550083X}
%
%
%%\bibitem{WE8compdata}
%%Shimada I.
%% {D}el {P}ezzo surfaces of degree one and examples of {Z}ariski
%%  multiples: computational data.
%% \url{http://www.math.sci.hiroshima-u.ac.jp/~shimada/K3andE. s.html},
%%  2025.
%
%\bibitem{ZariskiA}
%Zariski O.
% A theorem on the {P}oincar\'e{} group of an algebraic hypersurface.
%{Annals of Mathematics. Second Series}, 38(1):131--141, 1937.
% \url{https://doi.org/10.2307/1968515}
%
%\bibitem{ZariskiB}
%Zariski O.
% The topological discriminant group of a {R}iemann surface of genus
%  {$p$}.
%{American Journal of Mathematics}, 59(2):335--358, 1937.
% \url{https://doi.org/10.2307/2371416}

\end{thebibliography}



\end{document}
